\documentclass[UTF8, 12pt, fontset=fandol, oneside]{ctexart}
\usepackage{researchnotes}

\graphicspath{{tex/强化学习/figures/math-16/}}

\title{强化学习：神经科学}
\author{}
\date{}

\begin{document}

\maketitle

\section*{第 15 章\quad 神经科学}

神经科学是对神经系统的多学科研究：研究它们如何调节身体功能，如何控制行为，如何由于发育、学习和衰老而随时间改变，以及细胞和分子机制如何使这些功能成为可能。
强化学习最令人兴奋的方面之一，是来自神经科学的证据不断增加；这些证据表明，人类和许多其他动物的神经系统实现的算法，与强化学习算法有惊人的对应。
本章的主要目标是解释这些平行关系，以及它们对动物中与奖励有关的学习的神经基础有什么启示。

强化学习与神经科学之间最引人注目的接触点涉及多巴胺；多巴胺是一种化学物质，深度参与哺乳动物脑中的奖励加工。
多巴胺似乎把时间差分（TD）误差传递到发生学习和决策的脑结构。
这一平行关系由“多巴胺神经元活动的奖励预测误差假说”表达；这个假说来自计算强化学习与神经科学实验结果的汇合。
本章将讨论这个假说、导致它的神经科学发现，以及为什么它是理解脑奖励系统的一项重要贡献。
我们还会讨论强化学习与神经科学之间一些不如多巴胺/TD 误差平行关系那样醒目的对应；这些对应仍为思考动物中基于奖励的学习提供有用的概念工具。
强化学习的其他要素也有可能影响神经系统研究，但它们与神经科学的联系还相对不发达。
我们会讨论其中若干正在发展的联系，并且认为它们的重要性会随时间增长。

正如本书导论章的历史一节（第 1.7 节）所概述的，强化学习的许多方面都受到了神经科学影响。
本章的第二个目标，是让读者了解一些关于脑功能的思想；这些思想曾经促成我们处理强化学习的方式。
有些强化学习要素在脑功能理论的光照下更容易理解。
资格迹的思想尤其如此；资格迹是强化学习的基本机制之一，它最初来自对突触性质的猜想，而突触是神经细胞，也就是神经元，彼此通信所凭借的结构。

在本章中，我们不会很深入地探究动物中支撑基于奖励的学习的神经系统所具有的巨大复杂性：本章太短，而我们也不是神经科学家。
我们不会试图描述，甚至不会试图命名，许许多多被认为参与这些过程的脑结构和通路，也不会描述其中任何分子机制。
对于那些与强化学习也相当吻合的其他假说和模型，我们也无法充分给予公正介绍。
领域内专家有不同观点，这并不令人意外。
我们只能让读者瞥见这个迷人而正在发展的故事。
不过，我们希望本章能让你相信，一条非常富有成效的通道已经出现；它把强化学习及其理论基础与动物中基于奖励的学习的神经科学连接起来。

许多优秀出版物讨论强化学习与神经科学之间的联系，其中一些会在本章最后一节引用。
我们的处理与其中大多数不同，因为我们假定读者熟悉本书前面章节介绍的强化学习，但不假定读者了解神经科学。
我们先简要介绍理解后续内容所需的神经科学概念。

\subsection*{15.1\quad 神经科学基础}

为了跟上本章内容，了解一些关于神经系统的基本信息会很有帮助。
后面还会用到的术语在这里以斜体标出。
如果你已经有神经科学的初步知识，跳过本节也没有问题。

\emph{神经元}是神经系统的主要组成部分，是专门用电信号和化学信号处理并传递信息的细胞。
神经元有许多形式，但一个典型神经元具有胞体、\emph{树突}和一根\emph{轴突}。
树突是从胞体分支出来的结构，用来接收来自其他神经元的输入；对感觉神经元来说，它们也可以接收外部信号。
神经元的轴突是一根纤维，把该神经元的输出传送到其他神经元，或传送到肌肉和腺体。
神经元的输出由沿轴突传播的电脉冲序列组成，这些电脉冲称为\emph{动作电位}。
动作电位也称为\emph{脉冲}，当一个神经元产生脉冲时，就说它在\emph{放电}。
在神经网络模型中，常用实数表示神经元的\emph{放电率}，也就是某个时间单位内的平均脉冲数。

神经元的轴突可以广泛分支，使该神经元的动作电位到达许多目标。
神经元轴突的分支结构称为该神经元的\emph{轴突树}。
因为动作电位的传导是一个主动过程，有些像导火索的燃烧，所以当动作电位到达轴突分支点时，它会在所有传出分支上“点燃”动作电位；不过，向某个分支传播有时也会失败。
因此，具有大型轴突树的神经元的活动可以影响许多目标部位。

\emph{突触}是一种结构，通常位于轴突分支的末端，负责中介一个神经元到另一个神经元的通信。
突触把信息从\emph{突触前}神经元的轴突传递到\emph{突触后}神经元的树突或胞体。
除少数例外外，当动作电位到达突触前神经元时，突触会释放一种化学\emph{神经递质}。
这些例外是神经元之间直接电耦合的情形，不过这里不关心它们。
从突触前侧释放的神经递质分子扩散穿过\emph{突触间隙}，也就是突触前末梢与突触后神经元之间非常小的空间，然后与突触后神经元表面的受体结合，以兴奋或抑制其产生脉冲的活动，或以其他方式调节其行为。
一种特定神经递质可以结合到若干不同类型的受体，而每种受体对突触后神经元产生不同影响。
例如，多巴胺这种神经递质至少可以通过五种不同受体类型影响突触后神经元。
动物神经系统中已经鉴定出许多不同的化学物质作为神经递质。

神经元的\emph{背景}活动，是神经元似乎没有被实验者关心的任务相关突触输入所驱动时的活动水平，通常指它的放电率；例如，当神经元活动与实验中传递给受试对象的刺激不相关时。
背景活动可以是不规则的，原因可能来自更大网络的输入，也可能来自神经元或其突触内部的噪声。
有时，背景活动是神经元内在动态过程的结果。
与背景活动相反，神经元的\emph{时相}活动由突触输入通常引起的脉冲活动爆发构成。
无论是否作为背景活动，缓慢变化且常常呈等级式变化的活动，称为神经元的\emph{紧张性}活动。

在突触处释放的神经递质影响突触后神经元的强度或有效性，称为这个突触的\emph{效能}。
神经系统通过经验发生改变的一种方式，是突触效能发生改变；这些改变来自突触前神经元和突触后神经元活动的组合，有时也来自\emph{神经调质}的存在。
神经调质是一种神经递质，其作用不同于，或不止于，直接而快速的兴奋或抑制。

脑中包含若干不同的神经调制系统；每个系统由具有广泛分支轴突树的神经元簇组成，并且每个系统使用不同的神经递质。
神经调制可以改变神经回路的功能，中介动机、唤醒、注意、记忆、心境、情绪、睡眠和体温。
这里重要的是，神经调制系统可以分发某种类似标量信号的东西，例如强化信号，从而改变广泛分布的、对学习关键的部位中突触的运作。

突触效能发生改变的能力称为\emph{突触可塑性}。
它是负责学习的主要机制之一。
学习算法调节的参数或权重，对应于突触效能。
如下文详述，通过神经调质多巴胺来调制突触可塑性，是脑可能实现本书描述的许多学习算法的一种合理机制。

\subsection*{15.2\quad 奖励信号、强化信号、价值与预测误差}

神经科学与计算强化学习之间的联系，首先表现为脑中的信号与强化学习理论和算法中起突出作用的信号之间的平行关系。
在第 3 章中，我们说过，学习目标导向行为的任何问题，都可以归约为表示动作、状态和奖励的三个信号。
然而，为了解释神经科学与强化学习之间已经建立的联系，我们必须比这更不抽象，并考虑其他强化学习信号；这些信号在某些方面对应于脑中的信号。
除奖励信号外，这些信号还包括强化信号；我们将论证，强化信号不同于奖励信号；还包括价值信号，以及传达预测误差的信号。
当我们按某信号在强化学习理论中的功能给它贴标签时，就是在一个语境中这样做：在该语境里，这个信号对应于某个方程或算法中的一项。
另一方面，当我们说脑中的一个信号时，指的是一种生理事件，例如一次动作电位爆发或一种神经递质的分泌。
按神经信号的功能给它贴标签，例如把多巴胺神经元的时相活动称为强化信号，意味着这个神经信号的行为类似于，并且被猜想功能上类似于，相应的理论信号。

发现这些对应关系的证据涉及许多挑战。
与奖励加工有关的神经活动几乎可以在脑的每个部分发现，而且很难毫不含糊地解释结果，因为不同奖励相关信号的表征往往彼此高度相关。
实验需要被仔细设计，才能以某种确定程度把一种奖励相关信号与其他信号区分开，或把它与大量并不属于奖励加工的其他信号区分开。
尽管存在这些困难，许多实验仍以调和强化学习理论和算法的若干方面与神经信号为目标而展开，并且已经建立了一些令人信服的联系。
为了准备考察这些联系，本节剩余部分先提醒读者，各种奖励相关信号在强化学习理论中是什么意思。

在上一章末尾关于术语的说明中，我们说过，$R_t$ 类似于动物脑中的奖励信号，而不是动物环境中的对象或事件。
在强化学习中，奖励信号连同智能体环境定义了强化学习智能体试图解决的问题。
从这个意义上说，$R_t$ 类似于动物脑中的一个信号；它把初级奖励分配到遍布脑中的部位。
但是，动物脑中不太可能存在一个像 $R_t$ 这样的单一总控奖励信号。
最好把 $R_t$ 看作一种抽象，它概括了许多脑系统产生的大量神经信号的总体作用；这些系统评估感觉和状态的奖赏性质或惩罚性质。

强化学习中的\emph{强化信号}不同于奖励信号。
强化信号的功能是指导学习算法改变智能体的策略、价值估计或环境模型。
例如，对 TD 方法来说，时刻 $t$ 的强化信号是 TD 误差
\[
\delta_{t-1}=R_t+\gamma V(S_t)-V(S_{t-1}).
\]
\footnote{正如第 6.1 节提到的，在我们的记号中，$\delta_t$ 定义为 $R_{t+1}+\gamma V(S_{t+1})-V(S_t)$，所以 $\delta_t$ 直到时刻 $t+1$ 才可用。时刻 $t$ 可用的 TD 误差实际上是 $\delta_{t-1}=R_t+\gamma V(S_t)-V(S_{t-1})$。因为我们把时间步看得很小，甚至看作无穷小的时间间隔，所以不应把这个一步时间平移赋予过度重要性。}
某些算法的强化信号可以只是奖励信号，但对我们考虑的大多数算法来说，强化信号是由其他信息调节后的奖励信号，例如 TD 误差中的价值估计。

状态价值或动作价值的估计，也就是 $V$ 或 $Q$，说明什么东西从长期来看对智能体是好是坏。
它们是对智能体可以预期在未来累积的总奖励的预测。
智能体通过选择通向最大估计状态价值的状态的动作，或通过选择具有最大估计动作价值的动作，来作出好决策。

预测误差度量预期信号或感觉与实际信号或感觉之间的差异。
奖励预测误差（RPE）专门度量预期奖励信号和收到的奖励信号之间的差异；当奖励信号大于预期时为正，反之为负。
像式（6.5）这样的 TD 误差，是 RPE 的特殊种类；它们表示当前对长期奖励的预期与先前预期之间的差异。
当神经科学家谈到 RPE 时，他们通常，尽管并非总是，指 TD RPE；本章中我们简单地把它称为 TD 误差。
此外，在本章中，TD 误差通常指不依赖动作的 TD 误差；与之相对的是 Sarsa 和 Q-learning 等算法学习动作价值时使用的 TD 误差。
这是因为，与神经科学最著名的联系，是用不含动作的 TD 误差来陈述的；但我们并不想排除涉及依赖动作的 TD 误差的类似联系。
用于预测奖励以外信号的 TD 误差也很有用，不过这里不讨论这种情形；例如可见 Modayil、White 和 Sutton（2014）。

关于神经科学数据与这些理论定义信号之间的联系，可以提出许多问题。
一个观测到的信号更像奖励信号、价值信号、预测误差、强化信号，还是完全不同的东西？
如果它是误差信号，那么它是 RPE、TD 误差，还是类似 Rescorla--Wagner 误差（14.3）的更简单误差？
如果它是 TD 误差，它是否像 Q-learning 或 Sarsa 的 TD 误差那样依赖动作？
如上所示，探查脑以回答这类问题极其困难。
但是，实验证据提示，一种神经递质，也就是多巴胺，会发出 RPE 信号；进一步说，产生多巴胺的神经元的时相活动实际上传达 TD 误差（时相活动的定义见第 15.1 节）。
这一证据导致了多巴胺神经元活动的奖励预测误差假说，下面描述这一假说。

\subsection*{15.3\quad 奖励预测误差假说}

多巴胺神经元活动的\emph{奖励预测误差假说}提出，哺乳动物中产生多巴胺的神经元的时相活动，其功能之一，是把预期未来奖励的新旧估计之间的误差传递给遍布脑中的目标区域。
这个假说，虽然原文并非如此措辞，最早由 Montague、Dayan 和 Sejnowski（1996）明确提出；他们说明，强化学习中的 TD 误差概念可以解释多巴胺神经元时相活动的许多特征。
这些实验导致了这一假说；它们是在 20 世纪 80 年代和 90 年代早期，由神经科学家 Wolfram Schultz 的实验室完成的。
第 15.4 节描述这些有影响力的实验，第 15.6 节解释这些实验结果如何与 TD 误差相一致，而本章末尾的文献与历史说明一节包含一份围绕这一有影响力假说发展的文献指南。

Montague 等人（1996）把经典条件化 TD 模型的 TD 误差，与经典条件化实验中产生多巴胺的神经元的时相活动作了比较。
回忆第 14.2 节，经典条件化的 TD 模型基本上是带线性函数近似的半梯度下降 TD($\lambda$) 算法。
Montague 等人为建立这个比较作了若干假设。
第一，因为 TD 误差可以为负，而神经元不能有负放电率，所以他们假设，与多巴胺神经元活动对应的量是 $\delta_{t-1}+b_t$，其中 $b_t$ 是该神经元的背景放电率。
负 TD 误差对应于多巴胺神经元放电率降到其背景率以下。\footnote{在把 TD 误差与多巴胺神经元活动相联系的文献中，他们的 $\delta_t$ 与我们的 $\delta_{t-1}=R_t+\gamma V(S_t)-V(S_{t-1})$ 相同。}

第二个假设涉及每个经典条件化试次中访问的状态，以及这些状态如何表示为学习算法的输入。
这就是我们在第 14.2.4 节为 TD 模型讨论过的同一个问题。
Montague 等人选择了完整串行复合（CSC）表征，如图 14.1 左列所示；但在这里，短时长内部信号的序列会一直持续到 US 开始，而这里的 US 是非零奖励信号的到来。
这种表征允许 TD 误差模拟一个事实：多巴胺神经元活动不仅预测未来奖励，而且还对预测性线索之后奖励预期到来的\emph{时间}敏感。
必须有某种方式记录感官线索与奖励到达之间的时间。
如果一个刺激启动了一个在刺激结束后仍继续的内部信号序列，并且刺激后的每个时间步都有不同信号，那么刺激后的每个时间步都由一个不同状态表示。
因此，TD 误差由于依赖状态，就能够对一个试次内事件的时序敏感。

在采用这些关于背景放电率和输入表征的假设进行模拟试次时，TD 模型的 TD 误差与多巴胺神经元的时相活动惊人地相似。
预先概述下面第 15.4 节对这些相似性的细节描述，TD 误差与多巴胺神经元活动的以下特征平行：
1）多巴胺神经元的时相反应只在奖励事件未被预测时发生；
2）在学习早期，奖励之前的中性线索不会导致显著的时相多巴胺反应，但随着学习继续，这些线索获得预测价值，并开始诱发时相多巴胺反应；
3）如果一个更早的线索可靠地先于一个已经获得预测价值的线索，那么时相多巴胺反应会转移到更早的线索，并对较晚线索停止；
4）如果学习完成后，预测到的奖励事件被省略，那么多巴胺神经元的反应会在奖励事件预期发生时间之后不久降到基线水平以下。

尽管 Schultz 及其同事的实验中被监测的多巴胺神经元并不是每一个都以所有这些方式表现，但大多数被监测神经元的活动与 TD 误差之间的醒目对应，为奖励预测误差假说提供了强有力支持。
不过，也存在一些情形，其中基于该假说的预测与实验观察并不匹配。
输入表征的选择对于 TD 误差能多紧密地匹配多巴胺神经元活动的某些细节至关重要，尤其是关于多巴胺神经元反应时序的细节。
关于输入表征和 TD 学习其他特征，人们提出了不同想法，其中一些将在下文讨论，以使 TD 误差更好地拟合数据；不过，主要平行关系出现在 Montague 等人使用的 CSC 表征中。
总体而言，奖励预测误差假说已经被研究基于奖励的学习的神经科学家广泛接受，并且事实证明，它在神经科学实验结果不断积累的情况下仍格外有韧性。

为了准备描述支持奖励预测误差假说的神经科学实验，并提供一些背景，使读者能够理解该假说的重要性，我们接下来介绍关于多巴胺、它所影响的脑结构，以及它如何参与基于奖励的学习的一些已知内容。

\subsection*{15.4\quad 多巴胺}

多巴胺作为一种神经递质，由胞体主要位于哺乳动物中脑两簇神经元中的神经元产生：黑质致密部（SNpc）和腹侧被盖区（VTA）。
多巴胺在哺乳动物脑中的许多过程中起基本作用。
其中突出的是动机、学习、动作选择、大多数形式的成瘾，以及精神分裂症和帕金森病等障碍。
多巴胺称为神经调质，因为它执行许多不同于直接快速兴奋或抑制目标神经元的功能。
虽然关于多巴胺功能及其细胞作用细节还有许多未知，但显然，它是哺乳动物脑中奖励加工的基础。
多巴胺并不是参与奖励加工的唯一神经调质，而且它在厌恶情形，也就是惩罚中的作用仍有争议。
多巴胺在非哺乳动物中也可能以不同方式起作用。
但没有人怀疑，多巴胺对哺乳动物，包括人类，与奖励有关的过程是必不可少的。

一种早期的传统观点认为，多巴胺神经元向多个涉及学习和动机的脑区广播奖励信号。
这一观点来自 James Olds 和 Peter Milner 在 1954 年发表的一篇著名论文；该论文描述了电刺激大鼠脑中某些区域的作用。
他们发现，对特定区域的电刺激在控制大鼠行为方面充当极其强大的奖励：“……这种奖励对动物行为施加控制的程度极端，可能超过动物实验中先前使用过的任何其他奖励”（Olds 和 Milner，1954）。
后来的研究揭示，刺激最能产生这种奖赏效应的部位，会直接或间接兴奋通常由自然奖赏刺激兴奋的多巴胺通路。
在人类受试者身上也观察到类似作用。
这些观察强烈提示，多巴胺神经元活动发出奖励信号。

但是，如果奖励预测误差假说正确，即使它只解释多巴胺神经元活动的一部分特征，这种关于多巴胺神经元活动的传统观点也并不完全正确：多巴胺神经元的时相反应发出的是奖励预测误差信号，而不是奖励本身。
用强化学习术语说，多巴胺神经元在时刻 $t$ 的时相反应对应于
\[
\delta_{t-1}=R_t+\gamma V(S_t)-V(S_{t-1}),
\]
而不是 $R_t$。

强化学习理论和算法有助于把奖励预测误差观点与“多巴胺发出奖励信号”这一传统观念协调起来。
在本书讨论的许多算法中，$\delta$ 作为强化信号发挥作用，也就是说，它是学习的主要驱动力。
例如，$\delta$ 是经典条件化 TD 模型中的关键因素；在演员--评论家结构中，$\delta$ 是同时学习价值函数和策略的强化信号（第 13.5 节和第 15.7 节）。
依赖动作的 $\delta$ 形式是 Q-learning 和 Sarsa 的强化信号。
奖励信号 $R_t$ 是 $\delta_{t-1}$ 的关键组成部分，但在这些算法中，它并不完全决定 $\delta_{t-1}$ 的强化作用。
附加项 $\gamma V(S_t)-V(S_{t-1})$ 是 $\delta_{t-1}$ 的高阶强化部分；即使奖励发生（$R_t\neq 0$），如果奖励已经被完全预测，TD 误差也可能沉默；这一点将在下文第 15.6 节充分解释。

事实上，更仔细地看 Olds 和 Milner 1954 年的论文，会发现它主要讨论电刺激在工具性条件化任务中的强化作用。
电刺激不仅通过多巴胺对动机的影响激活动物行为，还会使大鼠很快学会按压杠杆来自我刺激；它们会在长时间内频繁这样做。
由电刺激触发的多巴胺神经元活动强化了大鼠的杠杆按压。

使用光遗传学方法的较新实验确证了多巴胺神经元时相反应作为强化信号的作用。
这些方法允许神经科学家在清醒、行为中的动物中，以毫秒尺度精确控制所选神经元类型的活动。
光遗传学方法把光敏蛋白引入所选神经元类型，使这些神经元能够借助激光闪光被激活或沉默。
第一个用光遗传学方法研究多巴胺神经元的实验表明，在小鼠中产生多巴胺神经元时相激活的光遗传刺激，足以使小鼠条件化地偏好箱室中收到这种刺激的一侧，而不是箱室中没有刺激或刺激频率更低的另一侧（Tsai 等，2009）。
另一个例子中，Steinberg 等人（2013）使用多巴胺神经元的光遗传激活，在奖励刺激被预期但被省略的时刻，在大鼠中制造人为的多巴胺神经元活动爆发；这些时刻通常会出现多巴胺神经元活动暂停。
当这些暂停被人为爆发替代时，通常会因缺少强化而下降的反应得以维持（在消退试次中），而通常会由于奖励已被预测而被阻滞的学习也得以发生（阻滞范式；第 14.2.1 节）。

关于多巴胺强化功能的额外证据来自果蝇的光遗传实验；不过在这些动物中，多巴胺的作用与其在哺乳动物中的作用相反：至少对被激活的多巴胺神经元群体而言，光触发的多巴胺神经元活动爆发在强化回避行为中，就像电足震一样发挥作用（Claridge-Chang 等，2009）。
虽然这些光遗传实验都没有表明多巴胺神经元时相活动具体像 TD 误差，但它们令人信服地证明，多巴胺神经元时相活动作为强化信号，在预测（经典条件化）和控制（工具性条件化）算法中都像 $\delta$ 那样发挥作用；在果蝇中，或许像 $-\delta$ 那样发挥作用。

多巴胺神经元特别适合向脑的许多区域广播强化信号。
这些神经元具有巨大的轴突树，每个神经元在 100 到 1000 倍于典型神经元轴突所达到的突触部位释放多巴胺。
右图显示的是单个多巴胺神经元的轴突树，它的胞体位于大鼠脑的 SNpc 中。
SNpc 或 VTA 多巴胺神经元的每根轴突，大约在目标脑区神经元的树突上形成 500000 个突触接触。

\begin{figure}[htbp]
\centering
\includegraphics[width=0.48\textwidth]{dopamine-axon-arbor.png}
\caption*{\textbf{图：}作为神经递质产生多巴胺的单个神经元的轴突树。这些轴突与目标脑区中大量神经元的树突形成突触接触。}
\end{figure}

如果多巴胺神经元广播像强化学习中的 $\delta$ 那样的强化信号，那么由于这是一个标量信号，也就是一个单一数字，SNpc 和 VTA 中的所有多巴胺神经元都应当或多或少相同地激活，使它们近乎同步地行动，把同一个信号发送到所有其轴突靶向的部位。
尽管人们通常认为多巴胺神经元确实以这种方式共同行动，现代证据却指向一个更复杂的图景：不同多巴胺神经元亚群会根据它们把信号发送到哪些结构，以及这些信号以哪些不同方式作用于其目标结构，而对输入作出不同反应。
多巴胺除了发出 RPE 信号还有其他功能；即使对确实发出 RPE 信号的多巴胺神经元来说，根据不同结构在产生强化行为中所扮演的角色，把不同 RPE 发送到不同结构也是有意义的。
这超出了本书详细处理的范围，但当决策可以分解为若干独立子决策时，或者更一般地说，作为处理信用分配问题的结构版本的一种方式，向量值 RPE 信号从强化学习角度看是有意义的：在可能参与产生某个决策的各组成结构之间，怎样分配成功的信用或失败的责任？
下文第 15.10 节还会对此略作说明。

大多数多巴胺神经元的轴突与额叶皮层和基底神经节中的神经元形成突触接触；这些脑区涉及自主运动、决策、学习，以及规划等认知功能。
因为把多巴胺与强化学习联系起来的大多数想法都关注基底神经节，并且来自多巴胺神经元的连接在那里特别密集，所以这里我们关注基底神经节。
基底神经节是一组位于前脑底部的神经元群，或称神经核。
基底神经节的主要输入结构称为纹状体。
基本上整个大脑皮层以及其他结构都向纹状体提供输入。
皮层神经元的活动传递关于感觉输入、内部状态和运动活动的大量信息。
皮层神经元的轴突与纹状体主要输入/输出神经元的树突形成突触接触；这些纹状体神经元称为中型棘状神经元。
来自纹状体的输出经由其他基底神经节核团和丘脑回环到皮层额叶区和运动区，使纹状体能够影响运动、抽象决策过程和奖励加工。
纹状体的两个主要分区对强化学习很重要：背侧纹状体主要涉及影响动作选择，而腹侧纹状体被认为对奖励加工的不同方面至关重要，包括把情感价值赋予感觉。

中型棘状神经元的树突布满棘突；皮层神经元的轴突在这些棘突的顶端形成突触接触。
同样与这些棘突形成突触接触的，是多巴胺神经元的轴突；在这种情况下，它们接触棘突茎（图 15.1）。
这种安排把皮层神经元的突触前活动、纹状体神经元的突触后活动和多巴胺聚到一起，使多种学习规则能够支配皮层--纹状体突触的可塑性。

\begin{figure}[htbp]
\centering
\includegraphics[width=0.72\textwidth]{figure-15-1.png}
\caption*{\textbf{图 15.1：}纹状体神经元的棘突，显示来自皮层神经元和多巴胺神经元的输入。皮层神经元轴突通过在纹状体神经元树突棘突顶端释放谷氨酸来影响纹状体神经元。VTA 或 SNpc 多巴胺神经元的轴突从棘突旁经过；这根轴突上的“多巴胺膨体”在棘突茎处或其附近释放多巴胺。}
\end{figure}

实际发生在这些棘突上的事情很复杂，而且尚未完全理解。
图 15.1 通过显示两种多巴胺受体、谷氨酸受体，以及各种信号能够相互作用的多种方式，暗示了这种复杂性；谷氨酸是皮层输入的神经递质。
不过，越来越多证据表明，从皮层到纹状体路径上的突触效能变化，也就是神经科学家所说的皮层--纹状体突触，关键依赖时机适当的多巴胺信号。

\subsection*{15.5\quad 奖励预测误差假说的实验证据}

多巴胺神经元会以活动爆发响应强烈、新颖或未预期的视觉和听觉刺激；这些刺激触发眼和身体运动，但这些神经元活动中很少有部分与运动本身有关。
这令人惊讶，因为多巴胺神经元退化是帕金森病的一个原因，而帕金森病症状包括运动障碍，尤其是自主发起运动的缺陷。
由于多巴胺神经元活动与刺激触发的眼和身体运动之间关系很弱，Romo 和 Schultz（1990）以及 Schultz 和 Romo（1990）迈出了走向奖励预测误差假说的第一步：他们在猴子移动手臂时记录多巴胺神经元活动和肌肉活动。

他们训练两只猴子，当猴子看到并听到一个小箱子的门打开时，从静止手位伸手进入箱子；箱中有一点苹果、一块饼干或一颗葡萄干。
猴子随后可以抓住食物并把它送入口中。
当猴子熟练之后，又在两个额外任务上训练。
第一个任务的目的是观察动作自发发起时多巴胺神经元会做什么。
箱子保持打开，但从上方遮盖，使猴子看不到里面，却能从下方伸手进去。
没有呈现触发刺激；当猴子伸手取到并吃下食物小块之后，实验者通常，尽管并非总是，会在猴子看不见也听不见的情况下，把食物粘到一根硬线上，悄悄把食物放回箱中。
在这里，Romo 和 Schultz 监测的多巴胺神经元活动也与猴子的动作无关；但是，当猴子第一次触到食物小块时，这些神经元中有很大比例产生时相反应。
当猴子只触到铁丝，或者在没有食物时探索箱子时，这些神经元没有反应。
这是很好的证据，说明这些神经元响应的是食物，而不是任务的其他方面。

Romo 和 Schultz 的第二个任务，目的是观察动作由刺激触发时会发生什么。
这个任务使用了另一个带可移动盖子的箱子。
箱盖打开的景象和声音触发猴子伸手到箱子里。
在这个情形中，Romo 和 Schultz 发现，经过一段训练后，多巴胺神经元不再响应触到食物，而是响应食物箱盖打开的景象和声音。
这些神经元的时相反应已经从奖励本身转移到预测奖励可得性的刺激上。
在一项后续研究中，Romo 和 Schultz 发现，他们监测的大多数多巴胺神经元，在行为任务语境之外，对箱盖打开的景象和声音没有反应。
这些观察提示，多巴胺神经元响应的既不是动作发起，也不是刺激的感官性质，而是在发出奖励预期信号。

Schultz 小组还进行了许多涉及 SNpc 和 VTA 多巴胺神经元的额外研究。
其中一系列实验特别有影响，因为它提示多巴胺神经元的时相反应对应于 TD 误差，而不是 Rescorla--Wagner 模型（14.3）中的误差。
在这些实验的第一个中（Ljungberg、Apicella 和 Schultz，1992），猴子被训练成在一个灯作为“触发线索”亮起后按下杠杆，以获得一滴苹果汁。
正如 Romo 和 Schultz 先前观察到的，许多多巴胺神经元最初响应奖励，也就是那滴果汁（图 15.2，上图）。
但是，随着训练继续，许多这类神经元失去了这种奖励反应，转而对预测奖励的灯光亮起作出反应（图 15.2，中图）。
随着训练继续，按杠杆变得更快，而响应触发线索的多巴胺神经元数量减少。

在这项研究之后，同一批猴子在一个新任务上训练（Schultz、Apicella 和 Ljungberg，1993）。
在这里，猴子面对两个杠杆，每个杠杆上方各有一盏灯。
其中一盏灯亮起，是一个“指令线索”，指示两个杠杆中的哪一个会产生一滴苹果汁。
在这个任务中，指令线索以固定的 1 秒间隔先于上一任务中的触发线索。
猴子学会在看到触发线索之前抑制伸手；多巴胺神经元活动增加，但现在，被监测多巴胺神经元的反应几乎完全发生在更早的指令线索，而不是触发线索上（图 15.2，下图）。
这里同样，当任务被学好时，响应指令线索的神经元数量大大减少。
在这些任务的学习过程中，多巴胺神经元活动从最初响应奖励，转移到响应更早的预测刺激，先推进到触发刺激，然后推进到更早的指令线索。
当反应在时间上向前移动时，它也从较晚刺激处消失。
这种反应向更早奖励预测物转移，同时丧失对较晚预测物的反应，是 TD 学习的一个标志性特征（例如见图 14.2）。

\begin{figure}[htbp]
\centering
\includegraphics[width=0.72\textwidth]{figure-15-2.png}
\caption*{\textbf{图 15.2：}多巴胺神经元的反应从对初级奖励的初始反应，转移到更早的预测刺激。图中绘出被监测多巴胺神经元在小时间间隔内产生的动作电位数，并对所有被监测多巴胺神经元取平均。上图：多巴胺神经元被未预期给出的一滴苹果汁激活。中图：随着学习，多巴胺神经元发展出对奖励预测性触发线索的反应，并失去对奖励给出的响应。下图：加入一个比触发线索早 1 秒的指令线索后，多巴胺神经元把反应从触发线索转移到更早的指令线索。}
\end{figure}

刚刚描述的任务揭示了多巴胺神经元活动与 TD 学习共有的另一个性质。
猴子有时会按错键，也就是按了不是被指令的那个键，于是没有得到奖励。
在这些试次中，许多多巴胺神经元在奖励通常给出的时间之后不久，放电率急剧降到基线以下；而且这在没有任何外部线索可标记通常奖励给出时间的情况下发生（图 15.3）。
猴子不知以某种方式在内部跟踪奖励的时机。
反应时序是最简单版本 TD 学习需要修改以解释多巴胺神经元反应某些时序细节的一个方面。
下一节将讨论这个问题。

\begin{figure}[htbp]
\centering
\includegraphics[width=0.55\textwidth]{figure-15-3.png}
\caption*{\textbf{图 15.3：}当预期奖励未发生时，多巴胺神经元反应会在奖励应当出现的时间之后不久降到基线以下。上图：多巴胺神经元被未预期给出的一滴苹果汁激活。中图：多巴胺神经元响应预测奖励的条件刺激（CS），而不响应奖励本身。下图：当 CS 预测的奖励未发生时，多巴胺神经元活动在奖励预期出现的时间之后不久降到基线以下。}
\end{figure}

上述研究中的观察使 Schultz 及其小组得出结论：多巴胺神经元响应未预测的奖励，响应最早的奖励预测物，而且如果奖励或奖励预测物没有在预期时间出现，多巴胺神经元活动会降到基线以下。
熟悉强化学习的研究者很快认识到，这些结果与 TD 误差在 TD 算法中作为强化信号时的行为极其相似。
下一节通过详细考察一个具体例子来探究这种相似性。

\subsection*{15.6\quad TD 误差/多巴胺对应}

本节解释 TD 误差 $\delta$ 与刚才描述的实验中观察到的多巴胺神经元时相反应之间的对应。
我们考察在类似上述任务的学习过程中，$\delta$ 如何变化；在该任务中，猴子先看到一个指令线索，然后在固定时间后必须对触发线索作出正确反应以获得奖励。
我们使用这个任务的一个简单理想化版本，但比通常更详细地讨论它，因为我们想强调 TD 误差与多巴胺神经元活动之间平行关系的理论基础。

第一个简化假设是，智能体已经学会了获得奖励所需的动作。
于是它的任务只是为它经历的状态序列学习对未来奖励的准确预测。
这就是一个预测任务，或者更技术地说，是一个策略评估任务：为固定策略学习价值函数（第 4.1 节和第 6.1 节）。
要学习的价值函数给每个状态赋予一个值，预测智能体按给定策略选择动作时该状态之后的回报；这里回报是所有未来奖励的和，可能带折扣。
作为猴子情境的模型，这并不现实，因为猴子很可能在学会正确行动的同时学习这些预测；强化学习算法如果既学习策略又学习价值函数，例如演员--评论家算法，也会如此。
不过，相比同时学习策略和价值函数的情形，这个场景更容易描述。

现在想象，智能体的经验分成多个试次；每个试次中，相同的状态序列重复出现，并且试次中的每个时间步出现一个不同状态。
再设想被预测的回报被限制为一个试次内的回报，这使试次类似于我们定义过的强化学习幕。
当然，实际被预测的回报并不局限于单个试次，而且试次之间的时间间隔也是决定动物学到什么的重要因素。
TD 学习同样如此，但这里我们假设回报不会跨多个试次累积。
在这个前提下，Schultz 及其同事所做实验中的一个试次，等同于强化学习的一幕；不过在本讨论中，为了更好地对应实验，我们使用“试次”而不是“幕”这个术语。

照常，我们还需要假设状态怎样被表示为学习算法的输入；这个假设会影响 TD 误差与多巴胺神经元活动对应得多紧密。
我们稍后讨论这个问题；目前假设使用 Montague 等人（1996）采用的同一个 CSC 表征，其中试次中每个时间步访问的每个状态都有一个分离的内部刺激。
这样，过程就化为本书第一部分讨论的表格情形。
最后，我们假设智能体用 TD(0) 学习价值函数 $V$，该函数存储在一个查找表中，并且所有状态初始化为零。
我们还假设这是一个确定性任务，折扣因子 $\gamma$ 非常接近 1，因此可以忽略。

图 15.4 显示了这个策略评估任务中，学习若干阶段的 $R$、$V$ 和 $\delta$ 的时间过程。
时间轴表示一个试次中状态序列被访问的时间间隔；为清楚起见，图中省略了各个状态。
奖励信号在每个试次中一直为零，直到智能体到达奖励状态；该状态显示在时间线右端附近，此时奖励信号变成某个正数，例如 $R^*$。
TD 学习的目标是预测一个试次中访问的每个状态的回报；在这个无折扣情形中，并且在我们假设预测只限于单个试次时，这个回报对每个状态都只是 $R^*$。

\begin{figure}[htbp]
\centering
\includegraphics[width=0.72\textwidth]{figure-15-4.png}
\caption*{\textbf{图 15.4：}TD 学习过程中 TD 误差 $\delta$ 的行为，与多巴胺神经元时相激活的特征一致。这里 $\delta$ 是时刻 $t$ 可用的 TD 误差，即 $\delta_{t-1}$。上：一串状态，即一段规则预测物，之后跟随非零奖励 $R^*$。学习早期：初始价值函数 $V$ 和初始 $\delta$；在开始时，$\delta$ 等于 $R^*$。学习完成：价值函数准确预测未来奖励，$\delta$ 在最早的预测状态为正，而在非零奖励时刻为 $0$。省略 $R^*$：在预测奖励被省略的时刻，$\delta$ 变为负。}
\end{figure}

在奖励状态之前，是一串奖励预测状态；其中最早的奖励预测状态显示在时间线的左端附近。
这类似于一个试次开始附近的状态，例如上文描述的 Schultz 等人（1993）猴子实验中由指令线索标记的状态。
它是一个试次中第一个可靠预测该试次奖励的状态。
当然，在现实中，前面试次中访问过的状态是更早的奖励预测状态；但因为我们把预测限制在单个试次内，它们不能算作\emph{这个}试次奖励的预测物。
下面我们会给出一个更令人满意但也更抽象的最早奖励预测状态描述。
一个试次中最晚的奖励预测状态，是紧接在该试次奖励状态之前的状态。
这是图 15.4 中时间线最右端附近的状态。
注意，一个试次的奖励状态并不预测该试次的回报：这个状态的价值会变成预测所有\emph{后续}试次的回报，而在这里的幕式表述中，我们假设该回报为零。

图 15.4 中标记为“学习早期”的图，显示了第一试次中 $V$ 和 $\delta$ 的时间过程。
因为除了到达奖励状态时之外，整个试次中的奖励信号都为零，而且所有 $V$ 值都为零，所以 TD 误差也为零，直到它在奖励状态处变成 $R^*$。
这是因为
\[
\delta_{t-1}=R_t+V_t-V_{t-1}=R_t+0-0=R_t,
\]
它一直为零，直到奖励发生时等于 $R^*$。
这里 $V_t$ 和 $V_{t-1}$ 分别是在一个试次中时刻 $t$ 和 $t-1$ 访问的状态的估计价值。
学习这一阶段的 TD 误差，类似于训练开始时多巴胺神经元对未预测奖励，例如一滴苹果汁，作出响应。

在这个第一试次以及所有后续试次中，TD(0) 更新发生在每一次状态转移上，如第 6 章所述。
这会逐步增加奖励预测状态的价值，增加从奖励状态向后传播，直到这些价值收敛到正确的回报预测。
在此情形中，由于我们假设没有折扣，所有奖励预测状态的正确预测都等于 $R^*$。
这可以在图 15.4 中标记为“学习完成”的 $V$ 图中看到；从最早到最晚奖励预测状态的所有状态价值都等于 $R^*$。
先于最早奖励预测状态的那些状态的价值保持较低，图 15.4 中显示为零，因为它们不是可靠的奖励预测物。

当学习完成，也就是 $V$ 达到正确值时，与从任意奖励预测状态出发的转移相关的 TD 误差为零，因为预测现在是准确的。
这是因为，对从一个奖励预测状态到另一个奖励预测状态的转移，有
\[
\delta_{t-1}=R_t+V_t-V_{t-1}=0+R^*-R^*=0,
\]
而对从最晚奖励预测状态到奖励状态的转移，有
\[
\delta_{t-1}=R_t+V_t-V_{t-1}=R^*+0-R^*=0.
\]
另一方面，从任意状态\emph{到}最早奖励预测状态的转移上的 TD 误差为正，因为该状态的低价值与随后奖励预测状态的较大价值之间存在不匹配。
事实上，如果先于最早奖励预测状态的状态价值为零，那么转移到最早奖励预测状态之后，有
\[
\delta_{t-1}=R_t+V_t-V_{t-1}=0+R^*-0=R^*.
\]
图 15.4 中“学习完成”的 $\delta$ 图显示了最早奖励预测状态处的这个正值，而其他各处为零。

转移到最早奖励预测状态时的正 TD 误差，类似于多巴胺反应持续出现在预测奖励的最早刺激上。
同样，当学习完成时，从最晚奖励预测状态到奖励状态的转移产生零 TD 误差，因为最晚奖励预测状态的价值由于正确而抵消了奖励。
这与很少有多巴胺神经元对完全预测的奖励而不是未预测奖励产生时相反应这一观察相平行。

学习之后，如果奖励突然被省略，TD 误差会在通常奖励时间变为负，因为最晚奖励预测状态的价值此时过高：
\[
\delta_{t-1}=R_t+V_t-V_{t-1}=0+0-R^*=-R^*,
\]
如图 15.4 中“省略 $R$”的 $\delta$ 图右端所示。
这类似于上文描述并见图 15.3 的 Schultz 等人（1993）实验中，当预期奖励被省略时，多巴胺神经元活动在相应时间降到基线以下。

最早奖励预测状态这一思想值得更多注意。
在上文描述的场景中，因为经验被分成试次，而且我们假设预测局限于单个试次，最早奖励预测状态总是一个试次的第一个状态。
显然，这是人为的。
一种更一般的思考方式是，把最早奖励预测状态理解为一个\emph{未被预测的奖励预测物}；这样的状态可以有很多个。
在动物一生中，许多不同状态都可能先于一个最早奖励预测状态。
然而，因为这些状态更经常后接其他并不预测奖励的状态，它们的奖励预测能力，也就是它们的价值，保持较低。
如果一个 TD 算法在动物整个生命中运行，它也会更新这些状态的价值；但这些更新不会一致地累积，因为按假设，没有任何这类状态可靠地先于一个最早奖励预测状态。
如果其中任何状态确实可靠地先于它，那么它们也会成为奖励预测状态。
这或许解释了为什么在过度训练下，即使是一个试次中最早奖励预测刺激上的多巴胺反应也会下降。
过度训练时，人们会预期，即使一个原先未被预测的预测物状态，也会被与更早状态相关的刺激所预测：动物在实验任务内外与环境的互动都会变得司空见惯。
然而，当引入一个新任务打破这种常规时，人们会看到 TD 误差重新出现；多巴胺神经元活动中确实观察到了这一点。

上文例子解释了为什么当动物在一个类似我们例子中理想化任务的任务中学习时，TD 误差会与多巴胺神经元的时相活动共享关键特征。
但多巴胺神经元时相活动并非每个性质都与 $\delta$ 的性质如此整齐地吻合。
最令人困扰的差异之一，涉及奖励比预期更早发生时会怎样。
我们已经看到，预期奖励的省略会在奖励预期时间产生负预测误差，对应于这种情况发生时多巴胺神经元活动降到基线以下。
如果奖励比预期更晚到达，那么它就是一个未预期奖励，并产生正预测误差。
TD 误差和多巴胺神经元反应都会发生这种情况。
但是，当奖励比预期更早到达时，多巴胺神经元并不做 TD 误差所做的事；至少在 Montague 等人（1996）和我们例子使用的 CSC 表征下是如此。
多巴胺神经元确实会响应早到奖励；这与正 TD 误差一致，因为奖励在那个时刻没有被预测。
然而，在较晚的奖励预期但被省略的时刻，TD 误差为负；与这一预测相反，多巴胺神经元活动并不按 TD 模型预测的方式降到基线以下（Hollerman 和 Schultz，1998）。
动物脑中发生的事情比带 CSC 表征的简单 TD 学习更复杂。

TD 误差与多巴胺神经元活动之间的一些不匹配，可以通过为 TD 算法选择合适参数值，以及使用 CSC 表征以外的刺激表征来处理。
例如，为处理刚描述的早到奖励不匹配，Suri 和 Schultz（1999）提出一种 CSC 表征，其中由较早刺激启动的内部信号序列会被奖励的发生取消。
Daw、Courville 和 Touretzky（2006）的另一项提议是，脑的 TD 系统使用由感觉皮层进行的统计建模产生的表征，而不是基于原始感觉输入的较简单表征。
Ludvig、Sutton 和 Kehoe（2008）发现，与使用 CSC 表征相比，带微刺激（MS）表征（图 14.1）的 TD 学习在早到奖励和其他情形中更好地拟合多巴胺神经元活动。
Pan、Schmidt、Wickens 和 Hyland（2005）发现，即使使用 CSC 表征，延长的资格迹也能改善 TD 误差对多巴胺神经元活动某些方面的拟合。
一般来说，TD 误差行为的许多精细细节依赖资格迹、折扣和刺激表征之间的微妙相互作用。
这类发现深化和细化了奖励预测误差假说，而没有驳倒其核心主张：多巴胺神经元的时相活动以发出 TD 误差信号为特征。

另一方面，TD 理论与实验数据之间还有一些差异，并不那么容易通过选择参数值和刺激表征来容纳；本章末尾文献与历史说明一节会提到其中一些差异。
随着神经科学家进行越来越精细的实验，更多不匹配很可能会被发现。
但奖励预测误差假说一直非常有效地充当一种催化剂，推动我们改进对脑奖励系统工作方式的理解。
复杂实验被设计出来，以验证或反驳从该假说导出的预测；反过来，实验结果又导致 TD 误差/多巴胺假说得到细化和展开。

这些发展有一个引人注目的方面：与多巴胺系统性质如此吻合的强化学习算法和理论，是从计算视角发展出来的，当时完全没有关于多巴胺神经元相关性质的任何知识。
要记得，TD 学习及其与最优控制和动态规划的联系，早在任何揭示多巴胺活动具有类 TD 性质的实验之前很多年就已经发展出来。
这种未规划的对应，尽管并不完美，却提示 TD 误差/多巴胺平行关系捕捉到了脑奖励过程中的某种重要东西。

除了解释多巴胺神经元时相活动的许多特征外，奖励预测误差假说还把神经科学连接到强化学习的其他方面，特别是连接到使用 TD 误差作为强化信号的学习算法。
神经科学距离完全理解多巴胺神经元时相活动的回路、分子机制和功能还很远，但支持奖励预测误差假说的证据，加上时相多巴胺反应是学习强化信号的证据，提示脑可能实现某种类似演员--评论家算法的东西，其中 TD 误差发挥关键作用。
其他强化学习算法当然也是合理候选，但演员--评论家算法与哺乳动物脑的解剖和生理特别吻合，如下两节所述。

\subsection*{15.7\quad 神经演员--评论家}

演员--评论家算法同时学习策略和价值函数。
“演员”是学习策略的成分，而“评论家”是学习当前由演员遵循的策略的一切有关内容、以“评论”演员动作选择的成分。
评论家使用 TD 算法来学习演员当前策略的状态价值函数。
价值函数使评论家能够通过向演员发送 TD 误差 $\delta$ 来评论演员的动作选择。
正 $\delta$ 表示该动作是“好的”，因为它导致一个价值高于预期的状态；负 $\delta$ 表示该动作是“坏的”，因为它导致一个价值低于预期的状态。
基于这些评论，演员不断更新其策略。

演员--评论家算法有两个独特特征，使人们认为脑可能实现某种类似算法。
第一，演员--评论家算法的两个成分，即演员和评论家，提示纹状体的两个部分，也就是背侧和腹侧分区（第 15.4 节），二者都对基于奖励的学习至关重要，可能分别起到类似演员和评论家的作用。
第二，演员--评论家算法中提示脑实现的另一个性质是，TD 误差作为演员和评论家的强化信号具有双重作用，尽管它对这两个成分中的学习影响不同。
这与神经回路的若干性质很好地吻合：多巴胺神经元的轴突同时靶向纹状体的背侧和腹侧分区；多巴胺似乎对调制这两个结构中的突触可塑性都至关重要；而像多巴胺这样的神经调质怎样作用于目标结构，取决于目标结构的性质，而不只取决于神经调质的性质。

第 13.5 节把演员--评论家算法作为策略梯度方法介绍，但 Barto、Sutton 和 Anderson（1983）的演员--评论家算法更简单，并且作为人工神经网络（ANN）给出。
这里我们描述一种类似 Barto 等人算法的 ANN 实现，并按照 Takahashi、Schoenbaum 和 Niv（2008），给出这种 ANN 如何由脑中的真实神经网络实现的示意性提议。
我们把演员和评论家的学习规则推迟到第 15.8 节讨论；在那里，我们把它们作为策略梯度形式的特例来给出，并讨论它们对多巴胺可能如何调制突触可塑性有什么启示。

图 15.5a 显示了一个把演员和评论家实现为成分网络的演员--评论家算法 ANN。
评论家由一个类似神经元的单元 $V$ 组成，其输出活动表示状态价值；还包括一个图中标为 TD 的菱形部件，它通过把 $V$ 的输出与奖励信号和先前状态价值结合起来计算 TD 误差，这也由从 TD 菱形回到它自身的回路所提示。
演员网络具有一层 $k$ 个演员单元，记为 $A_i$，$i=1,\ldots,k$。
每个演员单元的输出是一个 $k$ 维动作向量的一个分量。
另一种说法是，有 $k$ 个彼此竞争以被执行的分离动作，每个动作由一个演员单元发号施令；不过这里我们把整个 $A$ 向量看作一个动作。

\begin{figure}[htbp]
\centering
\includegraphics[width=0.92\textwidth]{figure-15-5.png}
\caption*{\textbf{图 15.5：}演员--评论家 ANN 及一种假想神经实现。a）作为 ANN 的演员--评论家算法。演员基于从评论家收到的 TD 误差 $\delta$ 调节策略；评论家用同一个 $\delta$ 调节状态价值参数。评论家由奖励信号 $R$ 和其状态价值估计的当前变化产生 TD 误差。b）演员--评论家算法的一种假想神经实现。演员和评论家中学习价值的部分分别放在纹状体背侧和腹侧分区。TD 误差由位于 VTA 和 SNpc 的多巴胺神经元传递，以调制从皮层区到背侧纹状体输入的突触效能变化。}
\end{figure}

评论家网络和演员网络都接收由多个特征组成的输入，这些特征表示智能体环境的状态。
回忆第 1 章，强化学习智能体的环境包括包含该智能体的“有机体”内外的成分。
图中把这些特征画成标为 $x_1,x_2,\ldots,x_n$ 的圆圈；为了使图保持简单，它们只画了两次。
从每个特征 $x_i$ 到评论家单元 $V$，以及到每个动作单元 $A_i$ 的每条连接，都有一个表示突触效能的权重。
评论家网络中的权重参数化价值函数，而演员网络中的权重参数化策略。
当这些权重按照下一节描述的评论家和演员学习规则改变时，网络就发生学习。

评论家中回路产生的 TD 误差，是改变评论家网络和演员网络中权重的强化信号。
在图 15.5a 中，这由标为“TD error $\delta$”的线表示；该线跨过评论家网络和演员网络中的所有连接。
网络实现的这一方面，连同奖励预测误差假说，以及多巴胺神经元活动被这些神经元广泛轴突树如此广泛分布这一事实，提示作为关于脑中与奖励有关的学习如何发生的假说，像这样的演员--评论家网络并非过分牵强。

图 15.5b 非常示意性地提示，根据 Takahashi 等人（2008）的假说，图左边的 ANN 可能怎样映射到脑结构。
该假说把演员和评论家中学习价值的部分分别放在纹状体的背侧和腹侧分区，也就是基底神经节的输入结构。
回忆第 15.4 节，背侧纹状体主要涉及影响动作选择，而腹侧纹状体被认为对奖励加工的不同方面至关重要，包括把情感价值赋予感觉。
大脑皮层连同其他结构向纹状体发送输入，传递关于刺激、内部状态和运动活动的信息。

在这种假想的演员--评论家脑实现中，腹侧纹状体把价值信息发送到 VTA 和 SNpc；在这些核团中，多巴胺神经元把它与奖励信息结合，以产生对应于 TD 误差的活动，不过多巴胺能神经元到底怎样计算这些误差还不清楚。
图 15.5a 中“TD error $\delta$”这条线，在图 15.5b 中变成标为“Dopamine”的线；它代表多巴胺神经元广泛分支的轴突，这些神经元的胞体位于 VTA 和 SNpc。
回看图 15.1，这些轴突与中型棘状神经元树突上的棘突形成突触接触；中型棘状神经元是纹状体背侧和腹侧分区的主要输入/输出神经元。
向纹状体发送输入的皮层神经元轴突在这些棘突的顶端形成突触接触。
按照该假说，正是在这些棘突处，从皮层区域到纹状体的突触效能变化，由关键依赖多巴胺提供的强化信号的学习规则支配。

图 15.5b 所示假说的一个重要含义是，多巴胺信号不是强化学习中标量 $R_t$ 那样的“总控”奖励信号。
事实上，该假说意味着，人们不一定能够探查脑，并在任何单个神经元的活动中记录到类似 $R_t$ 的信号。
许多互相连接的神经系统产生与奖励有关的信息，并且不同类型奖励会招募不同结构。
多巴胺神经元接收来自许多不同脑区的信息，因此图 15.5b 中标为“Reward”的 SNpc 和 VTA 输入，应理解为沿多个输入通道到达这些核团神经元的奖励相关信息向量。
理论上的标量奖励信号 $R_t$ 可能对应于所有奖励相关信息对多巴胺神经元活动的净贡献。
它是许多脑区中许多神经元活动模式的结果。

虽然图 15.5b 所示的演员--评论家神经实现可能在某些方面正确，但它显然需要细化、扩展和修改，才能够成为关于多巴胺神经元时相活动功能的成熟模型。
本章末尾的历史与文献说明一节引用了一些出版物，它们更详细地讨论这一假说的实验证据以及其不足之处。
现在我们详细考察演员和评论家学习算法对支配皮层--纹状体突触效能变化的规则有什么启示。

\subsection*{15.8\quad 演员和评论家的学习规则}

如果脑确实实现某种类似演员--评论家算法的东西，并且假设多巴胺神经元群向背侧和腹侧纹状体的皮层--纹状体突触广播一个共同强化信号，如图 15.5b 所示，尽管这很可能是上文提到的一种过度简化，那么这个强化信号会以不同方式影响这两个结构的突触。
评论家和演员的学习规则使用同一个强化信号，即 TD 误差 $\delta$，但它对这两个成分中学习的影响不同。
TD 误差与资格迹结合后，告诉演员如何更新动作概率，以到达价值更高的状态。
演员的学习类似于使用效果律型学习规则的工具性条件化（第 1.7 节）：演员努力使 $\delta$ 尽可能保持为正。
另一方面，TD 误差与资格迹结合后，告诉评论家应该以怎样的方向和大小改变价值函数参数，从而提高其预测准确性。
评论家使用类似经典条件化 TD 模型的学习规则（第 14.2 节），努力把 $\delta$ 的大小减小到尽可能接近零。
评论家和演员学习规则之间的差异很简单，但这个差异对学习有深刻影响，并且是演员--评论家算法工作方式的关键。
差异完全在于每种学习规则使用的资格迹。

在图 15.5b 这样的演员--评论家神经网络中，可以使用不止一组学习规则；为具体起见，这里关注第 13.6 节给出的带资格迹的持续性问题演员--评论家算法。
在每次从状态 $S_t$ 到 $S_{t+1}$ 的转移上，采取动作 $A_t$ 并收到奖励 $R_{t+1}$ 后，该算法计算 TD 误差 $\delta$，然后按如下方式更新资格迹向量（$\mathbf{z}_t^{\mathbf{w}}$ 和 $\mathbf{z}_t^\theta$）以及评论家和演员的参数（$\mathbf{w}$ 和 $\theta$）：
\[
\begin{aligned}
\delta_t &= R_{t+1}+\gamma \hat v(S_{t+1},\mathbf{w})-\hat v(S_t,\mathbf{w}),\\
\mathbf{z}_t^{\mathbf{w}} &= \lambda^{\mathbf{w}}\mathbf{z}_{t-1}^{\mathbf{w}}+\nabla \hat v(S_t,\mathbf{w}),\\
\mathbf{z}_t^\theta &= \lambda^\theta \mathbf{z}_{t-1}^\theta+\nabla \ln \pi(A_t|S_t,\theta),\\
\mathbf{w} &\leftarrow \mathbf{w}+\alpha^{\mathbf{w}}\delta_t\mathbf{z}_t^{\mathbf{w}},\\
\theta &\leftarrow \theta+\alpha^\theta\delta_t\mathbf{z}_t^\theta,
\end{aligned}
\]
其中 $\gamma\in[0,1)$ 是折扣率参数，$\lambda^{\mathbf{w}}\in[0,1]$ 和 $\lambda^\theta\in[0,1]$ 分别是评论家和演员的自举参数，而 $\alpha^{\mathbf{w}}>0$ 与 $\alpha^\theta>0$ 是类似的步长参数。

把近似价值函数 $\hat v$ 看成一个单一线性类神经元单元的输出；这个单元称为评论家单元，在图 15.5a 中标为 $V$。
于是，价值函数是状态 $s$ 的特征向量表示 $\mathbf{x}(s)=(x_1(s),\ldots,x_n(s))^\top$ 的线性函数，由权重向量 $\mathbf{w}=(w_1,\ldots,w_n)^\top$ 参数化：
\begin{equation}
\hat v(s,\mathbf{w})=\mathbf{w}^\top\mathbf{x}(s).
\tag{15.1}
\end{equation}
每个 $x_i(s)$ 都像是到达一个神经元突触的突触前信号，而该突触效能为 $w_i$。
按照上面的规则，评论家的权重增加 $\alpha^{\mathbf{w}}\delta_t\mathbf{z}_t^{\mathbf{w}}$；其中强化信号 $\delta_t$ 对应于广播到评论家单元所有突触的多巴胺信号。
评论家单元的资格迹向量 $\mathbf{z}_t^{\mathbf{w}}$，是 $\nabla\hat v(S_t,\mathbf{w})$ 的迹，也就是近期值的平均。
因为 $\hat v(s,\mathbf{w})$ 对权重是线性的，所以 $\nabla\hat v(S_t,\mathbf{w})=\mathbf{x}(S_t)$。

用神经术语说，这意味着每个突触都有自己的资格迹，它是向量 $\mathbf{z}_t^{\mathbf{w}}$ 的一个成分。
突触的资格迹按照到达该突触的活动水平累积，也就是按突触前活动水平累积；这里由到达该突触的特征向量 $\mathbf{x}(S_t)$ 的成分表示。
除此之外，该迹以由分数 $\lambda^{\mathbf{w}}$ 支配的速率向零衰减。
只要一个突触的资格迹非零，这个突触就\emph{有资格被修改}。
突触效能实际上如何被修改，取决于该突触仍有资格时到达的强化信号。
我们把这类评论家单元突触的资格迹称为\emph{非依存资格迹}，因为它们只依赖突触前活动，并不以任何方式依存于突触后活动。

评论家单元突触的非依存资格迹意味着，评论家单元的学习规则本质上就是第 14.2 节描述的经典条件化 TD 模型。
用上面给出的评论家单元及其学习规则的定义，图 15.5a 中的评论家与 Barto 等人（1983）的 ANN 演员--评论家中的评论家相同。
显然，像这样只由一个线性类神经元单元构成的评论家只是最简单起点；这个评论家单元是一个更复杂神经网络的代理，该网络能够学习复杂度更高的价值函数。

图 15.5a 中的演员是一个由 $k$ 个类神经元演员单元组成的单层网络，每个单元在时刻 $t$ 接收与评论家单元相同的特征向量 $\mathbf{x}(S_t)$。
每个演员单元 $j$，$j=1,\ldots,k$，有自己的权重向量 $\theta_j$；但由于所有演员单元相同，我们只描述其中一个单元并省略下标。
让这些单元遵循上面方程给出的演员--评论家算法的一种方式，是使每个单元成为一个 Bernoulli-logistic 单元。
这意味着每个演员单元在每个时刻的输出都是一个随机变量 $A_t$，取值为 0 或 1。
可以把取值 1 看成该神经元放电，也就是发出一个动作电位。
一个单元输入向量的加权和 $\theta^\top\mathbf{x}(S_t)$，通过指数 soft-max 分布（13.2）决定该单元的动作概率；对两个动作而言，这就是 logistic 函数：
\begin{equation}
\pi(1|s,\theta)=1-\pi(0|s,\theta)=\frac{1}{1+\exp(-\theta^\top\mathbf{x}(s))}.
\tag{15.2}
\end{equation}

如上所示，每个演员单元的权重按 $\theta\leftarrow\theta+\alpha^\theta\delta_t\mathbf{z}_t^\theta$ 增加，其中 $\delta$ 再次对应于多巴胺信号：发送到评论家单元所有突触的同一个强化信号。
图 15.5a 显示 $\delta_t$ 被广播到所有演员单元的所有突触，这使这个演员网络成为一个强化智能体团队；下文第 15.10 节会讨论这一点。
演员资格迹向量 $\mathbf{z}_t^\theta$ 是 $\nabla\ln\pi(A_t|S_t,\theta)$ 的迹，也就是近期值的平均。
为理解这个资格迹，可参见练习 13.5；该练习定义了这类单元，并要求给出其学习规则。
该练习要求通过计算梯度，把 $\nabla\ln\pi(a|s,\theta)$ 表示为 $a$、$\mathbf{x}(s)$ 和 $\pi(a|s,\theta)$ 的函数；对时刻 $t$ 实际发生的动作和状态，答案是
\begin{equation}
\nabla\ln\pi(A_t|S_t,\theta)=\big(A_t-\pi(A_t|S_t,\theta)\big)\mathbf{x}(S_t).
\tag{15.3}
\end{equation}

与评论家突触只累积突触前活动 $\mathbf{x}(S_t)$ 的非依存资格迹不同，演员单元突触的资格迹还依赖演员单元自身的活动。
我们把这称为\emph{依存资格迹}，因为它依存于这种突触后活动。
每个突触处的资格迹持续衰减，但会根据突触前神经元的活动以及突触后神经元是否放电而增加或减少。
式（15.3）中的因子 $A_t-\pi(A_t|S_t,\theta)$ 在 $A_t=1$ 时为正，否则为负。
\emph{演员单元资格迹中的突触后依存性，是评论家学习规则与演员学习规则之间的唯一区别。}
通过保存什么状态中采取了什么动作的信息，依存资格迹使奖励（正 $\delta$）的信用或惩罚（负 $\delta$）的责任，能够按照策略参数（演员单元突触效能）对这些单元输出的贡献，在这些参数之间分配；这些单元输出可能影响后来的 $\delta$ 值。
依存资格迹标记突触，使它们知道应当如何被修改，从而改变单元未来反应，使正 $\delta$ 的值更有可能出现。

评论家和演员学习规则对皮层--纹状体突触效能如何改变有什么启示？
两种学习规则都与 Donald Hebb 的经典提议有关：每当突触前信号参与激活突触后神经元时，该突触效能就增加（Hebb，1949）。
评论家和演员学习规则与 Hebb 提议共享这样一个思想：突触效能变化依赖若干因素的相互作用。
在评论家学习规则中，相互作用发生在强化信号 $\delta$ 与只依赖突触前信号的资格迹之间。
神经科学家把这称为\emph{二因子学习规则}，因为相互作用发生在两个信号或量之间。
另一方面，演员学习规则是\emph{三因子学习规则}，因为除了依赖 $\delta$ 外，它的资格迹还同时依赖突触前活动和突触后活动。
不过，与 Hebb 的提议不同，各因子的相对时机对突触效能如何改变至关重要；资格迹介入其中，使强化信号能够影响最近过去处于活动状态的突触。

关于演员和评论家学习规则的信号时序，还有一些细节值得更仔细注意。
在定义类神经元演员和评论家单元时，我们省略了突触输入影响真实神经元放电所需的一小段时间。
当来自突触前神经元的动作电位到达突触时，神经递质分子被释放，并扩散穿过突触间隙到达突触后神经元；在那里，它们与突触后神经元表面受体结合。
这会激活分子机制，使突触后神经元放电；或者在抑制性突触输入的情形中，抑制其放电。
这个过程可能需要几十毫秒。
不过，按照式（15.1）和式（15.2），评论家和演员单元的输入会瞬时产生该单元的输出。
在 Hebb 式可塑性的抽象模型中，忽略这种激活时间很常见；在这些模型中，突触效能按照同时发生的突触前活动和突触后活动的简单乘积而改变。
更现实的模型必须考虑激活时间。

对更现实的演员单元来说，激活时间尤其重要，因为它影响依存资格迹必须如何工作，才能把强化信用正确分配给合适的突触。
上面给出的演员单元学习规则中，定义依存资格迹的表达式
\[
\big(A_t-\pi(A_t|S_t,\theta)\big)\mathbf{x}(S_t)
\]
包括突触后因子 $\big(A_t-\pi(A_t|S_t,\theta)\big)$ 和突触前因子 $\mathbf{x}(S_t)$。
这之所以有效，是因为它忽略了激活时间，也就是忽略了突触前活动 $\mathbf{x}(S_t)$ 参与导致 $\big(A_t-\pi(A_t|S_t,\theta)\big)$ 中出现的突触后活动。
为了正确分配强化信用，定义资格迹的突触前因子必须是也定义该迹的突触后因子的原因。
更现实的演员单元的依存资格迹必须考虑激活时间。
激活时间不应与神经元接收到受该神经元活动影响的强化信号所需时间混淆。
资格迹的功能是跨越这个时间间隔，而它通常比激活时间长得多。
下一节将进一步讨论这一点。

神经科学对这一过程可能如何在脑中工作给出了一些线索。
神经科学家发现了一种称为\emph{脉冲时序依赖可塑性}（STDP）的 Hebb 式可塑性形式；它使脑中存在类似演员的突触可塑性更为可信。
STDP 是一种 Hebb 式可塑性，但突触效能变化依赖突触前动作电位与突触后动作电位的相对时序。
这种依赖可以有不同形式；但在研究最多的一种形式中，如果经该突触进入的脉冲在突触后神经元放电之前不久到达，则突触强度增加。
如果时序关系反过来，突触前脉冲在突触后神经元放电之后不久到达，则突触强度减小。
STDP 是一种考虑神经元激活时间的 Hebb 式可塑性，这是类似演员的学习所需成分之一。

STDP 的发现促使神经科学家研究一种三因子形式的 STDP 的可能性；在这种形式中，神经调制输入必须跟随时机适当的突触前和突触后脉冲。
这种突触可塑性称为\emph{奖励调制 STDP}，它非常类似于这里讨论的演员学习规则。
只有当突触前脉冲之后不久跟随突触后脉冲，并且在其后的某个时间窗内存在神经调制输入时，常规 STDP 本来会产生的突触变化才会发生。
越来越多证据表明，奖励调制 STDP 发生在背侧纹状体中型棘状神经元的棘突上，其中多巴胺提供神经调制因子；在图 15.5b 所示演员--评论家算法的假想神经实现中，这些部位正是演员学习发生的部位。
实验已经证明了奖励调制 STDP：只有当一个神经调制脉冲在一个时间窗内到达时，皮层--纹状体突触效能才会发生持久改变；这个时间窗可以持续到突触前脉冲之后紧随突触后脉冲后的 10 秒（Yagishita 等，2014）。
虽然证据是间接的，但这些实验指向具有较长时间进程的依存资格迹的存在。
产生这些迹的分子机制，以及可能支撑 STDP 的短得多的迹，目前还不清楚；不过，关于时间依赖和神经调质依赖突触可塑性的研究仍在继续。

我们在这里描述的类神经元演员单元，连同其效果律式学习规则，以稍简单的形式出现在 Barto 等人（1983）的演员--评论家网络中。
该网络受到生理学家 A. H. Klopf（1972，1982）提出的“享乐神经元”假说启发。
Klopf 假说的细节并不都与我们已经了解的突触可塑性一致，但 STDP 的发现以及奖励调制 STDP 证据的不断增加，提示 Klopf 的思想可能离目标并不远。
下面讨论 Klopf 的享乐神经元假说。

\subsection*{15.9\quad 享乐神经元}

在他的享乐神经元假说中，Klopf（1972，1982）猜想，单个神经元会通过调节其突触效能，来最大化被视为奖赏性的突触输入与被视为惩罚性的突触输入之间的差异；这种调节基于它自身动作电位带来的奖赏或惩罚后果。
换句话说，单个神经元可以用反应依存强化来训练，就像动物可以在工具性条件化任务中被训练一样。
他的假说包含这样一种思想：奖励和惩罚通过同一种突触输入传递给神经元；这种输入会兴奋或抑制该神经元产生脉冲的活动。
如果 Klopf 当时知道我们今天关于神经调制系统的知识，他或许会把强化作用赋给神经调制输入；但他想避免任何中心化的训练信息来源。
在 Klopf 的假说中，过去突触前活动和突触后活动的突触局部迹承担关键功能，使突触\emph{有资格}，这是他引入的术语，被后来的奖励或惩罚修改。
他猜想这些迹由每个突触局部的分子机制实现，因此不同于突触前神经元和突触后神经元的电活动。
在本章的文献与历史说明一节中，我们会提醒读者注意其他人提出过的一些类似方案。

Klopf 具体猜想，突触效能以下列方式改变。
当一个神经元发出一个动作电位时，它所有曾活跃并参与导致该动作电位的突触，都变得有资格经历效能变化。
如果该动作电位在适当时间段内后接奖励增加，那么所有有资格突触的效能增加。
对称地，如果该动作电位在适当时间段内后接惩罚增加，那么有资格突触的效能降低。
这是通过在突触前活动与突触后活动巧合时，或者更准确地说，在突触前活动与由该突触前活动参与导致的突触后活动配对时，在突触处触发资格迹来实现的；这就是我们所说的依存资格迹。
这本质上就是上一节描述的演员单元三因子学习规则。

Klopf 理论中资格迹的形状和时间进程，反映该神经元嵌入的许多反馈回路的持续时间。
其中一些反馈回路完全位于有机体的脑和身体内部，另一些则通过其运动和感觉系统，延伸到有机体的外部环境。
他的想法是，突触资格迹的形状类似于该神经元所嵌入反馈回路持续时间的直方图。
于是，资格迹的峰值会出现在该神经元参与最多的反馈回路的持续时间处。
本书描述的算法所使用的资格迹，是 Klopf 原始思想的简化版本，是由参数 $\lambda$ 和 $\gamma$ 控制的指数式，或几何式，递减函数。
这简化了模拟和理论，但我们把这些简单资格迹看作更接近 Klopf 原始构想的迹的占位物；后者会通过细化信用分配过程，在复杂强化学习系统中具有计算优势。

Klopf 的享乐神经元假说并不像乍看起来那样不可信。
一个研究充分的例子是单细胞寻找某些刺激并避开另一些刺激的情形，即细菌 \emph{Escherichia coli}。
这种单细胞生物的运动受其环境中化学刺激影响，这种行为称为趋化性。
它在液体环境中游动，方式是旋转附着在其表面的称为鞭毛的毛发状结构。
细菌环境中的分子会与其表面受体结合。
结合事件调制细菌反转鞭毛旋转的频率。
每次反转都会使细菌原地翻滚，然后朝随机的新方向离开。
少量化学记忆和计算会使鞭毛反转频率在细菌游向它生存所需分子，也就是吸引物，的较高浓度时降低；而在细菌游向有害分子，也就是排斥物，的较高浓度时增加。
结果是，细菌倾向于持续沿吸引物梯度向上游动，并倾向于避免沿排斥物梯度向上游动。

刚描述的趋化行为称为趋异运动。
它是一种试错行为，尽管其中不太可能涉及学习：细菌需要少量短时记忆来检测分子浓度梯度，但它大概不会保持长时记忆。
人工智能先驱 Oliver Selfridge 把这种策略称为“跑和摆动”，指出它作为基本适应性策略的用处：“如果情况在变好，就继续以同样方式前进；否则就到处移动”（Selfridge，1978，1984）。
类似地，人们可以把一个神经元想象成在由其所嵌入的复杂反馈回路集合组成的介质中“游动”，当然这不是字面意义上的游动；它行动以获得一种输入信号并避开其他输入信号。
不过，与细菌不同，神经元的突触强度会保留关于其过去试错行为的信息。
如果这种关于神经元，或只是某类神经元，行为的观点是可信的，那么该神经元与其环境交互的闭环性质对理解其行为很重要；其中，神经元的环境由动物的其余部分以及整个动物与之交互的环境共同构成。

Klopf 的享乐神经元假说超出了“单个神经元是强化学习智能体”这一思想。
他认为，智能行为的许多方面可以理解为大量自利的享乐神经元彼此交互的集体行为；它们构成动物神经系统这一巨大社会或经济系统。
无论这种神经系统观点是否有用，强化学习智能体的集体行为都对神经科学有启示。
下面讨论这一主题。

\subsection*{15.10\quad 集体强化学习}

强化学习智能体群体的行为，与社会和经济系统研究深度相关；而且如果类似 Klopf 享乐神经元假说的东西正确，它也与神经科学相关。
上文关于演员--评论家算法可能如何在脑中实现的假说，只狭窄地处理了这样一个事实的含义：纹状体背侧和腹侧分区，也就是按该假说分别作为演员和评论家位置的结构，各自包含数百万中型棘状神经元；这些神经元的突触经受由多巴胺神经元活动的时相爆发调制的变化。

图 15.5a 中的演员是一个由 $k$ 个演员单元组成的单层网络。
这个网络产生的动作是向量 $(A_1,A_2,\ldots,A_k)^\top$，假定它驱动动物行为。
所有这些单元的突触效能变化都依赖强化信号 $\delta$。
因为演员单元试图使 $\delta$ 尽可能大，$\delta$ 实际上对它们充当奖励信号；所以在这个情形中，强化与奖励相同。
因此，每个演员单元自身就是一个强化学习智能体；如果你愿意，也可以说是一个享乐神经元。
现在，为尽可能简化情形，假设每个这类单元在同一时间接收同一个奖励信号；不过，如上文指出的，假设多巴胺在同一条件和同一时刻释放到所有皮层--纹状体突触，很可能是过度简化。

强化学习理论能告诉我们，当一个强化学习智能体群体中的所有成员都按照共同奖励信号学习时会发生什么？
多智能体强化学习领域考虑强化学习智能体群体学习的许多方面。
虽然这个领域超出本书范围，但我们相信，其一些基本概念和结果与思考脑中弥散的神经调制系统有关。
在多智能体强化学习以及博弈论中，所有智能体都试图最大化它们同时收到的共同奖励信号的场景，称为协作博弈或团队问题。

团队问题有趣而困难之处在于，发送给每个智能体的共同奖励信号评价的是整个群体产生的活动\emph{模式}，也就是评价团队成员的\emph{集体动作}。
这意味着，任何单个智能体影响奖励信号的能力都很有限，因为任一单个智能体只贡献了共同奖励信号所评价的集体动作的一个成分。
在这种场景下，有效学习需要处理一个\emph{结构性信用分配问题}：哪些团队成员，或哪些团队成员群体，应当因有利奖励信号获得信用，或因不利奖励信号承担责任？
这是协作博弈，或团队问题，因为智能体们在寻求增加同一个奖励信号时是联合的：智能体之间没有利益冲突。
如果不同智能体接收不同奖励信号，而每个奖励信号同样评价群体的集体动作，并且每个智能体的目标都是增加自己的奖励信号，那么该场景就是竞争博弈。
在这种情况下，智能体之间可能存在利益冲突，意味着对某些智能体好的动作可能对另一些智能体坏。
甚至决定最佳集体动作应当是什么，也是博弈论中一个非平凡方面。
这种竞争设定也可能与神经科学相关，例如用于解释多巴胺神经元活动的异质性；但这里我们只关注协作或团队情形。

团队中的每个强化学习智能体怎样才能学会“做正确的事”，使团队的集体动作得到高度奖励？
一个有趣结果是，如果每个智能体都能有效学习，尽管它的奖励信号被大量噪声污染，而且它无法获得完整状态信息，那么整个群体会学会产生被共同奖励信号评价为改进的集体动作，即使智能体彼此不能通信。
每个智能体都面对自己的强化学习任务；在这个任务中，它对奖励信号的影响深深埋在其他智能体影响所产生的噪声中。
事实上，对任一智能体来说，所有其他智能体都是其环境的一部分，因为它的输入，无论是传递状态信息的部分还是奖励部分，都依赖所有其他智能体的行为方式。
此外，由于不能访问其他智能体的动作，甚至不能访问决定其策略的参数，每个智能体只能部分观察其环境状态。
这使每个团队成员的学习任务非常困难；但如果每个成员都使用一种即使在这些困难条件下也能增加奖励信号的强化学习算法，那么强化学习智能体团队可以学会产生随时间改进的集体动作，这种改进由团队的共同奖励信号评价。

如果团队成员是类神经元单元，那么每个单元都必须以随时间增加它收到的奖励量为目标，如第 15.8 节描述的演员单元那样。
每个单元的学习算法必须有两个基本特征。
第一，它必须使用依存资格迹。
回忆一下，用神经术语说，当某个突触的突触前输入参与导致突触后神经元放电时，该突触处会启动，或增加，一个依存资格迹。
相对地，非依存资格迹由突触前输入启动或增加，与突触后神经元做什么无关。
如第 15.8 节解释的，通过保存什么状态中采取了什么动作的信息，依存资格迹使奖励的信用或惩罚的责任，能够按照这些参数对决定智能体动作的贡献，分配给智能体策略参数。
通过类似推理，团队成员必须记住其近期动作，从而能够根据随后收到的奖励信号，提高或降低产生该动作的可能性。
依存资格迹的动作成分实现了这种动作记忆。
不过，由于学习任务的复杂性，依存资格只是信用分配过程的初步步骤：单个团队成员的动作与团队奖励信号变化之间的关系，是必须在许多试次中估计的统计相关。
在这一过程中，依存资格是必要但初步的一步。

带非依存资格迹的学习在团队设定中完全不起作用，因为它没有提供一种把动作与奖励信号后续变化相关起来的方式。
非依存资格迹足以学习预测，正如演员--评论家算法的评论家成分所做的那样；但它们不支持像演员成分必须做的那样学习控制。
一群类似评论家的智能体成员也许仍接收共同强化信号，但它们都会学习预测同一个量；在演员--评论家方法的情形中，这个量会是当前策略的期望回报。
群体每个成员在学习预测期望回报方面有多成功，将取决于它接收的信息，而不同群体成员收到的信息可能非常不同。
群体不需要产生差异化活动模式。
这不是这里定义的团队问题。

团队问题中集体学习的第二个要求，是团队成员的动作必须有变异性，以便团队探索集体动作空间。
强化学习智能体团队做到这一点的最简单方式，是让每个成员通过其输出的持续变异性独立探索自己的动作空间。
这会使整个团队改变其集体动作。
例如，第 15.8 节描述的演员单元团队会探索集体动作空间，因为每个单元作为 Bernoulli-logistic 单元，其输出概率性地依赖输入向量成分的加权和。
这个加权和使放电概率上升或下降，但总有变异性。
因为每个单元使用 REINFORCE 策略梯度算法（第 13 章），每个单元都会在随机探索其自身动作空间时，以最大化它经历的平均奖励率为目标调节自己的权重。
如 Williams（1992）所示，可以证明，一个 Bernoulli-logistic REINFORCE 单元团队整体上实现了关于团队共同奖励信号平均率的策略梯度算法，其中动作是团队的集体动作。

进一步地，Williams（1992）表明，当团队中的单元相互连接形成多层 ANN 时，使用 REINFORCE 的 Bernoulli-logistic 单元团队会沿平均奖励梯度上升。
在这种情况下，奖励信号被广播到所有单元
的网络中，尽管奖励可能只依赖网络输出单元的集体动作。
这意味着，一个由 Bernoulli-logistic REINFORCE 单元组成的多层团队，像一个用广泛使用的误差反向传播方法训练的多层网络那样学习；但在这里，反向传播过程被广播的奖励信号替代。
在实践中，误差反向传播方法要快得多，但强化学习团队方法作为神经机制更可信，尤其是考虑到第 15.8 节讨论的关于奖励调制 STDP 的知识。

团队成员通过独立探索来进行探索，只是团队探索的最简单方式。
如果团队成员通过彼此通信或响应共同输入来协调动作，从而聚焦于集体动作空间中的特定部分，那么更复杂的方法也是可能的。
也有比依存资格迹更复杂的机制来处理结构性信用分配；当可能的集体动作集合以某种方式受到限制时，这在团队问题中更容易。
一个极端情形是赢者通吃安排，例如脑中侧抑制的结果；它把集体动作限制为只有一个或少数团队成员贡献的动作。
在这种情形下，赢者获得结果奖励或惩罚的信用或责任。

协作博弈，或团队问题，以及非协作博弈问题中的学习细节，超出本书范围。
本章末尾的文献与历史说明一节引用了一组选定的相关出版物，其中包括关于集体强化学习对神经科学启示的研究的广泛参考。

\subsection*{15.11\quad 脑中的基于模型方法}

强化学习中无模型算法与基于模型算法之间的区分，正在证明有助于思考动物学习和决策过程。
第 14.6 节讨论了这种区分如何与动物行为中习惯性行为和目标导向行为之间的区分相一致。
上文讨论的脑可能如何实现演员--评论家算法的假说，只与动物的习惯性行为模式有关，因为基本的演员--评论家方法是无模型的。
哪些神经机制负责产生目标导向行为，它们又如何与习惯性行为背后的机制相互作用？

研究这些行为模式所涉及脑结构的一种方式，是使大鼠脑的某个区域失活，然后在结果贬值实验（第 14.6 节）中观察大鼠做什么。
这类实验的结果表明，上文描述的演员--评论家假说把演员放在背侧纹状体中，过于简单。
使背侧纹状体的一部分，也就是背外侧纹状体（DLS），失活，会损害习惯学习，使动物更多依赖目标导向过程。
另一方面，使背内侧纹状体（DMS）失活，会损害目标导向过程，迫使动物更多依赖习惯学习。
这类结果支持这样一种观点：啮齿动物中的 DLS 更多参与无模型过程，而它们的 DMS 更多参与基于模型过程。
在人类受试者和非人灵长类中，用功能神经成像进行的类似实验研究结果，也支持这样一种观点：灵长类脑中的类似结构以不同方式参与习惯性和目标导向行为模式。

其他研究则识别出与基于模型过程相关的活动，它们位于人脑前额叶皮层；前额叶皮层是额叶最前部，涉及包括规划和决策在内的执行功能。
具体涉及的是眶额皮层（OFC），即前额叶皮层中紧邻眼睛上方的部分。
人类功能神经成像，以及猴子中单个神经元活动记录，揭示 OFC 中有强活动与具有生物意义的刺激的主观奖励价值有关，也有活动与动作结果所预期的奖励有关。
尽管并非没有争议，这些结果提示 OFC 显著参与目标导向选择。
它可能对动物环境模型中的奖励部分至关重要。

参与基于模型行为的另一个结构是海马体；它是对记忆和空间导航至关重要的结构。
大鼠海马体在大鼠以目标导向方式穿越迷宫的能力中起关键作用；这种能力使 Tolman 提出动物在选择动作时使用模型或认知地图的思想（第 14.5 节）。
海马体也可能是人类想象新经验能力的关键组成部分（Hassabis 和 Maguire，2007；Ólafsdóttir、Barry、Saleem、Hassabis 和 Spiers，2015）。

最直接把海马体牵涉到规划中的发现，来自解码海马体中神经元活动的实验；规划是使用环境模型作决策所需的过程。
这些实验确定海马活动在逐时刻基础上表示空间的哪一部分。
当大鼠在迷宫中的选择点暂停时，海马体中的空间表征会沿着动物从该点可能采取的路径向前扫过，而不是向后扫过（Johnson 和 Redish，2007）。
此外，这些扫过所表示的空间轨迹与大鼠随后导航行为密切对应（Pfeiffer 和 Foster，2013）。
这些结果提示，海马体对动物环境模型的状态转移部分至关重要，而且它是一个系统的一部分；该系统使用模型来模拟可能的未来状态序列，以评估可能行动路线的后果，也就是一种规划形式。

上述结果加入了一个庞大文献，这些文献讨论目标导向，或基于模型，学习和决策的底层神经机制；但许多问题仍未回答。
例如，像 DLS 和 DMS 这样结构上相似的区域，怎样能够成为无模型算法和基于模型算法这样不同的学习与行为模式的基本组成部分？
是否有分离的结构分别负责环境模型中我们所说的转移成分和奖励成分？
所有规划是否都在决策时通过可能未来行动路线的模拟来进行，就像海马体中的前向扫过活动所提示的那样？
换言之，所有规划是否都类似第 8.10 节的展开算法？
或者，模型是否有时在后台参与，以细化或重新计算价值信息，如 Dyna 架构（第 8.2 节）所示？
脑如何在习惯系统和目标导向系统之间仲裁？
事实上，这些系统的神经基质之间是否存在清晰分离？

证据并不指向最后一个问题的肯定答案。
Doll、Simon 和 Daw（2012）总结情况时写道：“基于模型的影响似乎或多或少无处不在，只要脑加工奖励信息”，即使在被认为对无模型学习至关重要的区域中也是如此。
这包括多巴胺信号本身；除了作为无模型过程基础的奖励预测误差外，多巴胺信号也可能表现出基于模型信息的影响。

由强化学习的无模型/基于模型区分所启发的持续神经科学研究，有潜力使我们更清楚地理解脑中的习惯性过程和目标导向过程。
更好地掌握这些神经机制，可能导致算法以计算强化学习尚未探索的方式结合无模型方法和基于模型方法。

\subsection*{15.12\quad 成瘾}

理解药物滥用的神经基础，是神经科学中的高优先级目标；它有可能为这一严重公共健康问题产生新治疗。
一种观点是，药物渴求来自同一套动机和学习过程；这些过程使我们寻求服务于生物需要的自然奖励经验。
成瘾性物质由于具有强强化性，有效劫持了我们的自然学习和决策机制。
这种观点有其合理性，因为许多，尽管不是全部，滥用药物会直接或间接提高纹状体中多巴胺神经元轴突末端周围区域的多巴胺水平；纹状体是一个与正常基于奖励的学习牢固相关的脑结构（第 15.7 节）。
但是，与药物成瘾相关的自毁行为并不是正常学习的特征。
当奖励来自成瘾药物时，多巴胺中介学习有什么不同？
成瘾是否是正常学习对一些物质作出响应的结果，而这些物质在我们的进化史中基本不可获得，因此进化无法选择抵御其损害作用？
或者，成瘾性物质是否以某种方式干扰正常的多巴胺中介学习？

多巴胺神经元活动的奖励预测误差假说及其与 TD 学习的联系，是 Redish（2004）关于成瘾某些特征模型的基础；当然，这个模型并不能解释所有特征。
该模型基于这样一种观察：施用可卡因和某些其他成瘾药物，会产生短暂的多巴胺增加。
在该模型中，假设这种多巴胺激增以一种不能由价值函数变化抵消的方式增加 TD 误差 $\delta$。
换言之，正常奖励被先前事件预测到时，$\delta$ 会相应降低（第 15.6 节）；但当奖励信号变得被预测时，由成瘾刺激对 $\delta$ 的贡献并不会降低：药物奖励不能被“预测掉”。
该模型通过阻止奖励信号来自成瘾药物时 $\delta$ 变为负，从而做到这一点；这样就消除了与给药相关状态的 TD 学习误差校正特征。
结果是这些状态的价值无界增加，使导致这些状态的动作优先于所有其他动作。

成瘾行为比 Redish 模型给出的这个结果复杂得多，但该模型的主要思想可能是拼图的一块。
或者该模型也可能具有误导性。
多巴胺似乎并不在所有形式成瘾中都起关键作用，而且并不是每个人都同样容易发展成瘾行为。
此外，该模型没有纳入慢性用药伴随的许多回路和脑区变化，例如导致药物效应随反复使用而减弱的变化。
成瘾也很可能涉及基于模型的过程。
尽管如此，Redish 模型说明了强化学习理论如何被用于理解一个重大健康问题。
以类似方式，强化学习理论已经影响了计算精神病学这一新领域的发展；该领域旨在通过数学和计算方法改进对精神障碍的理解。

\subsection*{15.13\quad 小结}

脑奖励系统涉及的神经通路复杂且尚未完全理解，但旨在理解这些通路及其在行为中作用的神经科学研究正在快速推进。
这项研究正在揭示脑奖励系统与本书所介绍强化学习理论之间的醒目对应。

多巴胺神经元活动的\emph{奖励预测误差假说}，由一些科学家提出；他们认识到 TD 误差行为与产生多巴胺这种神经递质的神经元活动之间存在醒目的平行关系。
多巴胺对哺乳动物中与奖励有关的学习和行为至关重要。
20 世纪 80 年代晚期和 90 年代，在神经科学家 Wolfram Schultz 的实验室中进行的实验证明，多巴胺神经元只有在动物不预期奖励事件时，才会以称为时相反应的大量活动爆发响应这些事件；这提示多巴胺神经元发出的是奖励预测误差信号，而不是奖励本身。
进一步地，这些实验显示，随着动物学会基于先前感觉线索预测奖励事件，多巴胺神经元的时相活动会转移到更早的预测线索，同时对较晚预测线索的活动会下降。
这与强化学习智能体学会预测奖励时，TD 误差的后备效应相平行。

其他实验结果牢固地确立了：多巴胺神经元的时相活动是学习的强化信号，它借助产生多巴胺的神经元高度分支的轴突到达脑的多个区域。
这些结果与我们在奖励信号 $R_t$ 和强化信号之间作出的区分一致；在我们介绍的大多数算法中，强化信号是 TD 误差 $\delta_t$。
多巴胺神经元的时相反应是强化信号，而不是奖励信号。

一个突出假说是，脑实现某种类似演员--评论家算法的东西。
脑中的两个结构，即纹状体背侧和腹侧分区，都在基于奖励的学习中起关键作用；它们可能分别像演员和评论家那样发挥作用。
TD 误差同时作为演员和评论家的强化信号这一点，与以下事实很好地吻合：多巴胺神经元轴突同时靶向纹状体背侧和腹侧分区；多巴胺似乎对调制这两个结构中的突触可塑性至关重要；而且像多巴胺这样的神经调质对目标结构的影响，取决于目标结构的性质，而不只取决于神经调质的性质。

演员和评论家可以由 ANN 实现，这些 ANN 由类神经元单元组成，并具有基于第 13.5 节所述策略梯度演员--评论家方法的学习规则。
这些网络中的每条连接都类似于脑中神经元之间的突触，而学习规则对应于支配突触效能如何改变的规则；这种改变是突触前和突触后神经元活动，以及对应于多巴胺神经元输入的神经调制输入共同作用的函数。
在这种设定中，每个突触都有自己的资格迹，记录涉及该突触的过去活动。
演员和评论家学习规则之间的唯一区别，是它们使用不同种类的资格迹：评论家单元的迹是非依存的，因为它们不涉及评论家单元输出；而演员单元的迹是依存的，因为除演员单元输入外，它们还依赖演员单元输出。
在脑中演员--评论家系统的假想实现中，这些学习规则分别对应于支配皮层--纹状体突触可塑性的规则；这些突触把信号从皮层传到纹状体背侧和腹侧分区的主要神经元，同时也接收来自多巴胺神经元的输入。

演员--评论家网络中演员单元的学习规则，与\emph{奖励调制脉冲时序依赖可塑性}密切对应。
在脉冲时序依赖可塑性（STDP）中，突触前和突触后活动的相对时序决定突触变化方向。
在奖励调制 STDP 中，突触变化还依赖一种神经调质，例如多巴胺，在满足 STDP 条件之后持续可达 10 秒的时间窗内到达。
越来越多证据表明，奖励调制 STDP 发生在皮层--纹状体突触处；在演员--评论家系统的假想神经实现中，这里正是演员学习发生的地方，这提高了某种类似演员--评论家系统存在于某些动物脑中的假说的可信度。

突触资格和演员学习规则基本特征的思想，来自 Klopf 的“享乐神经元”假说（Klopf，1972，1981）。
他猜想，单个神经元会根据自身动作电位的奖赏性或惩罚性后果来调节其突触效能，从而寻求获得奖励并避免惩罚。
神经元活动可以影响其后来的输入，因为神经元嵌入许多反馈回路之中；其中一些在动物神经系统和身体内部，另一些则穿过动物的外部环境。
Klopf 关于资格的思想是，如果突触参与了神经元放电，那么突触会被暂时标记为有资格修改；这使它成为依存形式的资格迹。
如果在突触有资格时到达强化信号，该突触的效能就被修改。
我们提到细菌的趋化行为，作为单个细胞指引其运动以寻找某些分子并避开另一些分子的例子。

多巴胺系统的一个显著特征是，释放多巴胺的纤维广泛投射到脑的多个部分。
虽然很可能只有某些多巴胺神经元群体广播同一个强化信号，但如果这个信号到达许多参与演员型学习的神经元突触，那么这种情形
可以被建模为一个\emph{团队问题}。
在这类问题中，强化学习智能体集合中的每个智能体都接收同一个强化信号，而该信号依赖集合或团队中所有成员的活动。
如果每个团队成员都使用足够有能力的学习算法，那么即使团队成员彼此之间没有直接通信，团队也能集体学会改进整个团队的表现；这种表现由全局广播的强化信号评价。
这与脑中多巴胺信号的广泛分散一致，并为训练多层网络的广泛使用的误差反向传播方法提供了一种神经上可信的替代。

无模型强化学习与基于模型强化学习之间的区分，正在帮助神经科学家研究习惯性学习和目标导向学习与决策的神经基础。
到目前为止，研究指出某些脑区更多参与一种过程而不是另一种过程，但图景仍不清楚，因为无模型过程和基于模型过程似乎并没有在脑中整齐分离。
许多问题仍未回答。
也许最令人感兴趣的是，有证据表明海马体，这一传统上与空间导航和记忆相关的结构，似乎参与模拟可能的未来行动路线，作为动物决策过程的一部分。
这提示它是使用环境模型进行规划的系统的一部分。

强化学习理论也正在影响人们关于药物滥用背后神经过程的思考。
一个关于药物成瘾某些特征的模型，基于奖励预测误差假说。
该模型提出，成瘾刺激，例如可卡因，会使 TD 学习不稳定，从而导致与服药有关的动作价值无界增长。
这距离完整成瘾模型还很远，但它说明计算视角如何提出可通过进一步研究检验的理论。
新的计算精神病学领域也类似地关注使用计算模型，其中一些来自强化学习，以更好地理解精神障碍。

本章只是触及了强化学习神经科学与计算机科学和工程中强化学习发展之间如何相互影响的表面。
强化学习算法的大多数特征，其设计都来自纯粹计算考虑；但也有一些特征受到关于神经学习机制的假说影响。
值得注意的是，随着关于脑奖励过程的实验数据不断积累，强化学习算法中许多纯粹由计算动机产生的特征，正在显示出与神经科学数据一致。
计算强化学习的其他特征，例如资格迹，以及强化学习智能体团队在全局广播强化信号影响下学会集体行动的能力，也可能随着神经科学家继续揭示基于奖励的动物学习和行为的神经基础，而最终显示出与实验数据平行。

\section*{文献与历史说明}

处理学习和决策神经科学与本书介绍的强化学习方法之间平行关系的出版物数量巨大。
我们只能引用很小一部分。
Niv（2009）、Dayan 和 Niv（2008）、Glimcher（2011）、Ludvig、Bellemare 和 Pearson（2011），以及 Shah（2012）都是不错的起点。

与经济学、进化生物学和数理心理学一起，强化学习理论正在帮助形成关于人类和非人灵长类选择神经机制的定量模型。
由于本章关注学习，所以只是轻轻触及决策神经科学。
Glimcher（2003）引入了“神经经济学”领域；在该领域中，强化学习从经济学视角为研究决策的神经基础作出贡献。
另见 Glimcher 和 Fehr（2013）。
Dayan 和 Abbott（2001）关于神经科学中的计算和数学建模的教材，包含了强化学习在这些方法中的作用。
Sterling 和 Laughlin（2015）从使高效适应性行为成为可能的一般设计原则出发，考察了学习的神经基础。

\noindent\textbf{15.1}
基础神经科学有许多好的说明。
Kandel、Schwartz、Jessell、Siegelbaum 和 Hudspeth（2013）是权威且非常全面的资料来源。

\noindent\textbf{15.2}
Berridge 和 Kringelbach（2008）综述了奖励和愉悦的神经基础，指出奖励加工具有许多维度并涉及许多神经系统。
篇幅不允许讨论 Berridge 和 Robinson（1998）有影响力的研究；他们区分了刺激的享乐影响，也就是所谓“喜欢”，与动机效应，也就是所谓“想要”。
Hare、O'Doherty、Camerer、Schultz 和 Rangel（2008）从经济学视角考察了价值相关信号的神经基础，区分目标价值、决策价值和预测误差。
决策价值是目标价值减去动作成本。
另见 Rangel、Camerer 和 Montague（2008）、Rangel 和 Hare（2010），以及 Peters 和 Büchel（2010）。

\noindent\textbf{15.3}
多巴胺神经元活动的奖励预测误差假说，最突出地由 Schultz、Dayan 和 Montague（1997）讨论。
这一假说最早由 Montague、Dayan 和 Sejnowski（1996）明确提出。
按他们表述的假说，它指奖励预测误差（RPE），但并未专门指 TD 误差；不过，他们对假说的发展清楚表明，他们指的是 TD 误差。
我们所知最早认识到 TD 误差/多巴胺联系的，是 Montague、Dayan、Nowlan、Pouget 和 Sejnowski（1993）；他们提出了一种由 Schultz 小组多巴胺信号结果启发的 TD 误差调制 Hebb 学习规则。
Quartz、Dayan、Montague 和 Sejnowski（1992）的摘要也指出了这一联系。
Montague 和 Sejnowski（1994）强调脑中预测的重要性，并概述了受 TD 误差调制的预测性 Hebb 学习，如何通过一种弥散神经调制系统，例如多巴胺系统，实现。
Friston、Tononi、Reeke、Sporns 和 Edelman（1994）提出了脑中价值依赖学习的一个模型，其中突触变化由全局神经调制信号提供的类 TD 误差中介，尽管他们没有特别指出多巴胺。
Montague、Dayan、Person 和 Sejnowski（1995）提出了蜜蜂采蜜使用 TD 误差的模型。
该模型基于 Hammer、Menzel 及同事的研究（Hammer 和 Menzel，1995；Hammer，1997），这些研究显示神经调质章鱼胺在蜜蜂中作为强化信号。
Montague 等人（1995）指出，多巴胺很可能在脊椎动物脑中扮演类似角色。
Barto（1995a）把演员--评论家结构与基底神经节回路联系起来，并讨论了 TD 学习与 Schultz 小组主要结果之间的关系。
Houk、Adams 和 Barto（1995）提出，TD 学习和演员--评论家结构可能如何映射到基底神经节的解剖、生理和分子机制。
Doya 和 Sejnowski（1998）扩展了他们早先关于鸟鸣学习模型的论文（Doya 和 Sejnowski，1995），把一种等同于多巴胺的类 TD 误差纳入其中，用来强化对要记忆的听觉输入的选择。
O'Reilly 和 Frank（2006）以及 O'Reilly、Frank、Hazy 和 Watz（2007）主张，时相多巴胺信号是 RPE，但不是 TD 误差。
为支持其理论，他们引用了可变刺激间隔的结果；这些结果不匹配简单 TD 模型的预测；也引用了高阶条件化超过二阶条件化很少被观察到，而 TD 学习并不如此受限这一观察。
Dayan 和 Niv（2008）讨论了强化学习理论和奖励预测误差假说如何与实验数据相一致的“好、坏、丑”。
Glimcher（2011）综述了支持奖励预测误差假说的经验发现，并强调该假说对当代神经科学的重要意义。

\noindent\textbf{15.4}
Graybiel（2000）是关于基底神经节的简短入门。
文中提到的涉及多巴胺神经元光遗传激活的实验，由 Tsai、Zhang、Adamantidis、Stuber、Bonci、de Lecea 和 Deisseroth（2009），Steinberg、Keiflin、Boivin、Witten、Deisseroth 和 Janak（2013），以及 Claridge-Chang、Roorda、Vrontou、Sjülson、Li、Hirsh 和 Miesenböck（2009）完成。
Fiorillo、Yun 和 Song（2013）、Lammel、Lim 和 Malenka（2014），以及 Saddoris、Cacciapaglia、Wightman 和 Carelli（2015）等研究显示，多巴胺神经元的信号性质针对不同目标区域而专门化。
RPE 信号神经元可能属于具有不同目标并服务不同功能的多个多巴胺神经元群体之一。
Eshel、Tian、Bukwich 和 Uchida（2016）发现，在小鼠经典条件化期间，外侧 VTA 中多巴胺神经元的奖励预测误差反应具有同质性；不过他们的结果并不排除更广区域中反应的多样性。
Gershman、Pesaran 和 Daw（2009）研究了可以分解成具有分离奖励信号的独立成分的强化学习任务，并在人类神经成像数据中发现证据，提示脑利用这种结构。

\noindent\textbf{15.5}
Schultz 1998 年的综述文章，是进入关于多巴胺神经元奖励预测信号的极其庞大文献的良好入口。
Berns、McClure、Pagnoni 和 Montague（2001），Breiter、Aharon、Kahneman、Dale 和 Shizgal（2001），Pagnoni、Zink、Montague 和 Berns（2002），以及 O'Doherty、Dayan、Friston、Critchley 和 Dolan（2003）描述了支持人脑中存在类 TD 误差信号的功能脑成像研究。

\noindent\textbf{15.6}
本节在解释 TD 误差如何模拟 Schultz 小组关于多巴胺神经元时相反应的主要结果时，大体遵循 Barto（1995a）。

\noindent\textbf{15.7}
本节主要基于 Takahashi、Schoenbaum 和 Niv（2008）以及 Niv（2009）。
据我们所知，Barto（1995a）和 Houk、Adams 与 Barto（1995）最早推测演员--评论家算法在基底神经节中的可能实现。
基于人类受试者在工具性条件化中进行功能磁共振成像，O'Doherty、Dayan、Schultz、Deichmann、Friston 和 Dolan（2004）提示，演员和评论家最可能分别位于背侧纹状体和腹侧纹状体。
Gershman、Moustafa 和 Ludvig（2014）关注基底神经节强化学习模型中时间如何表示，讨论了各种时间表示计算方法的证据和含义。

本节描述的演员--评论家结构的假想神经实现，只包含已知基底神经节解剖和生理的很少细节。
除了 Houk、Adams 和 Barto（1995）的更详细假说之外，许多其他假说包含与解剖和生理更具体的联系，并声称可以解释更多数据。
这些包括 Suri 和 Schultz（1998，1999）、Brown、Bullock 和 Grossberg（1999）、Contreras-Vidal 和 Schultz（1999）、Suri、Bargas 和 Arbib（2001）、O'Reilly 和 Frank（2006），以及 O'Reilly、Frank、Hazy 和 Watz（2007）提出的假说。
Joel、Niv 和 Ruppin（2002）批判性评价了其中若干模型在解剖上的可信性，并提出一个替代方案，以容纳基底神经节回路中一些被忽视的特征。

\noindent\textbf{15.8}
这里讨论的演员学习规则比 Barto 等人（1983）早期演员--评论家网络中的规则更复杂。
该网络中的演员单元资格迹只是 $A_t\times\mathbf{x}(S_t)$ 的迹，而不是完整的 $(A_t-\pi(A_t|S_t,\theta))\mathbf{x}(S_t)$。
那项工作没有受益于第 13 章介绍的策略梯度理论，也没有受益于 Williams（1986，1992）的贡献；后者说明 Bernoulli-logistic 单元 ANN 如何实现策略梯度方法。

Reynolds 和 Wickens（2002）提出了一种皮层--纹状体通路中突触可塑性的三因子规则，其中多巴胺调制皮层--纹状体突触效能的变化。
他们讨论了这种学习规则的实验证据及其可能分子基础。
脉冲时序依赖可塑性（STDP）的决定性展示归功于 Markram、Lübke、Frotscher 和 Sakmann（1997）；Levy 和 Steward（1983）等人的较早实验也表明，突触前和突触后脉冲的相对时序对于诱发突触效能变化至关重要。
Rao 和 Sejnowski（2001）提出，STDP 可能是具有约 10 毫秒持续时间非依存资格迹的突触处类 TD 机制的结果。
Dayan（2002）评论说，这需要 Sutton 和 Barto（1981a）早期经典条件化模型中的误差，而不是一个真正的 TD 误差。
奖励调制 STDP 庞大文献中的代表性出版物包括 Wickens（1990）、Reynolds 和 Wickens（2002），以及 Calabresi、Picconi、Tozzi 和 Di Filippo（2007）。
Pawlak 和 Kerr（2008）表明，多巴胺对于在中型棘状神经元的皮层--纹状体突触处诱发 STDP 是必要的。
另见 Pawlak、Wickens、Kirkwood 和 Kerr（2010）。
Yagishita、Hayashi-Takagi、Ellis-Davies、Urakubo、Ishii 和 Kasai（2014）发现，只有在 STDP 刺激后 0.3 到 2 秒的时间窗内，多巴胺才促进小鼠中型棘状神经元的棘突增大。
Izhikevich（2007）提出并探索了使用 STDP 时序条件来触发依存资格迹的思想。
Frémaux、Sprekeler 和 Gerstner（2010）提出了基于奖励调制 STDP 的规则成功学习的理论条件。

\noindent\textbf{15.9}
Klopf 的享乐神经元假说（Klopf，1972，1982）启发了我们作为 ANN 实现的演员--评论家算法；该 ANN 使用单个类神经元单元，即演员单元，来实现类似效果律的学习规则（Barto、Sutton 和 Anderson，1983）。
与 Klopf 的突触局部资格有关的思想也由其他人提出。
Crow（1968）提出，皮层神经元突触的变化对神经活动后果敏感。
他强调，在奖励调制形式的突触可塑性中，需要处理神经活动与其后果之间的时间延迟；他提出一种依存形式的资格，但与整个神经元而不是单个突触相关。
按照他的假说，一波神经元活动会导致参与该波的细胞发生短时变化，使它们从未如此激活的细胞背景中被挑选出来；这些细胞因这种短时变化而对奖励信号敏感；如果这种信号在变化衰减时间结束之前发生，那么细胞之间的突触连接会变得更有效（Crow，1968）。

Crow 反对先前关于回响神经回路发挥这一作用的提议；他指出，奖励信号对这种回路的作用会“建立导致回响的突触连接（也就是说，那些在奖励信号时活动中涉及的连接），而不是通向适应性运动输出的路径上的连接”。
Crow 进一步假定，奖励信号通过一种“独特神经纤维系统”传递；大概就是 Olds 和 Milner（1954）曾触及的那个系统；它会把突触连接“从短时形式转化为长时形式”。

在另一个很有远见的假说中，Miller 提出了一种类似效果律的学习规则，其中包含突触局部的依存资格迹：
\begin{quote}
设想在一个特定感觉情境中，神经元 B 偶然发出一次“有意义的活动爆发”，随后这种活动爆发被转化为运动动作，而运动动作又改变情境。必须假设，这个有意义爆发在神经元层面上影响其自身所有在当时活跃的突触，从而对将被加强的突触作出初步选择，尽管还没有真正加强它们。加强信号随后作出最终选择，并在合适突触中建立确定性改变。（Miller，1981，第 81 页）
\end{quote}
Miller 的假说还包含一种类似评论家的机制；他称之为“感觉分析器单元”。
该机制按照经典条件化原则工作，向神经元提供强化信号，使它们学会从较低价值状态移动到较高价值状态，从而预见了在演员--评论家结构中使用 TD 误差作为强化信号。
Miller 的思想不仅与 Klopf 的思想平行，除了它明确调用一个不同的“加强信号”之外，而且也预见了奖励调制 STDP 的一般特征。

Seung（2003）称为“享乐突触”的一个相关但不同思想，是突触以效果律方式各自调节其释放神经递质的概率：如果释放之后有奖励，那么释放概率增加；如果未释放之后有奖励，那么释放概率降低。
这本质上与 Minsky 在其 1954 年 Princeton 博士论文中使用的学习方案相同；在那里，他把类似突触的学习元件称为 SNARC（随机神经模拟强化计算器）。
这些思想也涉及个体突触的依存资格，不过它依存于单个突触的活动，而不是突触后神经元的活动。
Unnikrishnan 和 Venugopal（1994）的方案也相关；该方案使用 Harth 和 Tzanakou（1974）的基于相关的方法来调节 ANN 权重。

Frey 和 Morris（1997）提出“突触标签”的思想，用于诱发突触效能的长时程增强。
尽管这与 Klopf 的资格并不完全不同，但这种标签被假设为一个突触的临时增强；随后神经元激活可以把它转化为持久增强。
O'Reilly 和 Frank（2006）以及 O'Reilly、Frank、Hazy 和 Watz（2007）的模型使用工作记忆来跨越时间间隔，而不是使用资格迹。
Wickens 和 Kotter（1995）讨论了突触资格的可能机制。
He、Huertas、Hong、Tie、Hell、Shouval 和 Kirkwood（2015）提供证据，支持皮层神经元突触中存在依存资格迹；其时间进程类似 Klopf 假定的资格迹。

把神经元使用与细菌趋化性相关学习规则的比喻，由 Barto（1989）讨论。
Koshland 对细菌趋化性的广泛研究，在一定程度上受到了细菌特征与神经元特征之间相似性的启发（Koshland，1980）。
另见 Berg（1975）。
Shimansky（2009）提出了一种与上述 Seung 方案有些相似的突触学习规则；在该规则中，每个突触各自像趋化细菌那样行动。
在这种情况下，一个突触集合会在突触权重值的高维空间中“游向”吸引物。
Montague、Dayan、Person 和 Sejnowski（1995）提出了一个蜜蜂采蜜行为的趋化性式模型，其中涉及神经调质章鱼胺。

\noindent\textbf{15.10}
强化学习智能体在团队问题和博弈问题中的行为研究有很长历史，大致经历了三个阶段。
据我们所知，第一阶段始于俄罗斯数学家和物理学家 M. L. Tsetlin 的研究。
他于 1966 年去世后，其工作集以 Tsetlin（1973）出版。
我们的第 1.7 节和第 4.8 节在多臂老虎机问题相关处提到了他对学习自动机的研究。
Tsetlin 文集还包括团队问题和博弈问题中学习自动机的研究；这些研究导致了该领域后来的工作，使用随机学习自动机，如 Narendra 和 Thathachar（1974，1989）、Viswanathan 和 Narendra（1974）、Lakshmivarahan 和 Narendra（1982）、Narendra 和 Wheeler（1983），以及 Thathachar 和 Sastry（2002）所述。
Thathachar 和 Sastry（2011）是较新的综合性说明。
这些研究大多局限于非关联学习自动机，意思是它们没有处理关联式，或语境式，多臂老虎机问题（第 2.9 节）。

第二阶段始于把学习自动机扩展到关联式，或语境式，情形。
Barto、Sutton 和 Brouwer（1981）以及 Barto 和 Sutton（1981b）在单层 ANN 中实验了关联随机学习自动机，其中全局强化信号被广播。
该学习算法是 Harth 和 Tzanakou（1974）的 Alopex 算法的一种关联式扩展。
Barto 等人把实现这种学习的类神经元元件称为关联搜索元件（ASE）。
Barto 和 Anandan（1985）引入了一种称为关联奖励--惩罚（$A_{R-P}$）算法的关联强化学习算法。
他们通过把随机学习自动机理论与模式分类理论相结合，证明了一个收敛结果。
Barto（1985，1986）以及 Barto 和 Jordan（1987）描述了 $A_{R-P}$ 单元团队连接成多层 ANN 的结果，显示它们可以在全局广播强化信号下学习非线性函数，例如 XOR 等。
Barto（1985）广泛讨论了这种 ANN 方法，以及这类学习规则与当时文献中其他规则的关系。
Williams（1992）数学分析并拓宽了这类学习规则，并把它们的使用与用于训练多层 ANN 的误差反向传播方法联系起来。
Williams（1988）描述了若干将反向传播与强化学习结合起来训练 ANN 的方式。
Williams（1992）表明，$A_{R-P}$ 算法的一个特例是 REINFORCE 算法，尽管一般 $A_{R-P}$ 算法取得了更好结果（Barto，1985）。

第三阶段对强化学习智能体团队的兴趣，受到以下因素影响：人们越来越理解多巴胺作为广泛广播神经调质的作用，并推测奖励调制 STDP 的存在。
与早期研究相比，这些研究更多考虑突触可塑性的细节以及来自神经科学的其他约束。
相关出版物按时间和字母顺序包括：Bartlett 和 Baxter（1999，2000）、Xie 和 Seung（2004）、Baras 和 Meir（2007）、Farries 和 Fairhall（2007）、Florian（2007）、Izhikevich（2007）、Pecevski、Maass 和 Legenstein（2008）、Legenstein、Pecevski 和 Maass（2008）、Kolodziejski、Porr 和 Wörgötter（2009）、Urbanczik 和 Senn（2009），以及 Vasilaki、Frémaux、Urbanczik、Senn 和 Gerstner（2009）。
Nowé、Vrancx 和 De Hauwere（2012）综述了更广泛多智能体强化学习领域的较新发展。

\noindent\textbf{15.11}
Yin 和 Knowlton（2006）综述了啮齿动物结果贬值实验的发现；这些发现支持这样一种观点：习惯性行为和目标导向行为，按心理学家的用语，分别最常与背外侧纹状体（DLS）和背内侧纹状体（DMS）中的加工相关。
Valentin、Dickinson 和 O'Doherty（2007）在人类受试者结果贬值设定中进行功能成像实验，其结果提示眶额皮层（OFC）是目标导向选择的重要成分。
Padoa-Schioppa 和 Assad（2006）在猴子中记录单元活动，支持 OFC 在编码指导选择行为的价值中发挥作用。
Rangel、Camerer 和 Montague（2008）以及 Rangel 和 Hare（2010）从神经经济学视角综述了关于脑如何作出目标导向决策的发现。
Pezzulo、van der Meer、Lansink 和 Pennartz（2014）综述了内部生成序列的神经科学，并提出这些机制如何可能成为基于模型规划成分的模型。
Daw 和 Shohamy（2008）提出，虽然多巴胺信号与习惯性，或无模型，行为很好地连接，但目标导向，或基于模型，行为还涉及其他过程。
Bromberg-Martin、Matsumoto、Hong 和 Hikosaka（2010）的实验数据表明，多巴胺信号包含与习惯性行为和目标导向行为都相关的信息。
Doll、Simon 和 Daw（2012）认为，脑中服务于习惯性学习与选择和目标导向学习与选择的机制之间，可能并没有清晰分离。

\noindent\textbf{15.12}
Keiflin 和 Janak（2015）综述了 TD 误差与成瘾之间的联系。
Nutt、Lingford-Hughes、Erritzoe 和 Stokes（2015）批判性评价了成瘾是由多巴胺系统障碍导致的假说。
Montague、Dolan、Friston 和 Dayan（2012）概述了计算精神病学领域的目标和早期努力，而 Adams、Huys 和 Roiser（2015）综述了较新进展。

\end{document}
